\chapter{Value distribution}\label{ch:oc-values}

The distribution of $\log\zeta(\sigma+it)$. Off the critical line this is the Bohr--Jessen law, for each fixed $\sigma>\frac12$, with no hypothesis on the zeros: the Euler product converges and the proof uses only the mean square of a short mollifier. On the line it is Selberg's central limit theorem. The argument moves a distance $W/\log T$ off the line, with $W$ a small power of $\log\log\log T$, and compares $\log|\zeta|$ with a short prime sum; it does not use an unproved zero-density estimate. Every appeal to the Riemann hypothesis in the exercises is written in the exercise. The mean square of $|\zeta|^2$ on the critical line is a different theorem, proved in Volume~14, and is not this one. The covariance of the prime sums at two fixed heights is Exercise~\ref{ex:twopoint} below. A multipoint theorem for shifts that shrink with $T$ is not proved in this volume.

\section{The Bohr--Jessen law}\label{ch:bohrjessen}

\begin{center}\itshape
Walk up a vertical line to the right of $\frac12$ and record the values of $\zeta$.\\
They wander over the whole plane, with a law that you can compute,\\
because the primes, seen from far away, are independent.
\end{center}

Few analytic number theorists own an Olympic medal. Harald Bohr does: he played for the Danish football team that won silver at the London Olympics of 1908, two years before his doctoral thesis. His thesis was about Dirichlet series, and in the years that followed, with Richard Courant and Edmund Landau, he studied the values of $\zeta(\sigma+it)$ on a vertical line $\sigma>\frac12$, which nobody before him had thought of as a single object. In 1914 Bohr and Courant proved that for $\frac12<\sigma\le1$ these values are dense in the complex plane: $\zeta$ comes as close as you like to every complex number. Figure~\ref{fig:zetacurve} shows the curve $t\mapsto\zeta(0.7+it)$ beginning its wandering.

The idea behind Bohr's results is the theme of this book in its purest form. The Euler product $\zeta(s)=\prod_p(1-p^{-s})^{-1}$ depends on $t$ only through the numbers $p^{-it}=e^{-it\log p}$, points on the unit circle. Because the numbers $\log p$ are linearly independent over $\mathbb Q$ (that is unique factorisation), these points behave, for a random $t$, like \emph{independent} uniform points: the Steinhaus random multiplicative function of Chapter~\ref{ch:randommult}. In 1930 and 1932 Bohr and Børge Jessen turned this into a limit theorem: $\zeta(\sigma+it)$, for $t$ uniform in a long interval, converges in law to a random Euler product. This chapter proves it, deduces the density of the values, and shows that the limit law of $\log|\zeta|$ has a smooth density. On the critical line itself the random Euler product diverges, and Section~\ref{ch:selbergclt} takes over.

Throughout, $\frac12<\sigma<1$ is fixed, $t$ is uniform on $[T,2T]$, and $\mathbb E_T$ and $P_T$ refer to this choice. $(X_p)_p$ are independent and uniform on the unit circle, and for $|x|=1$,
\[
\ell_p(x)=-\log\Bigl(1-\frac x{p^\sigma}\Bigr)=\sum_{k\ge1}\frac{x^k}{kp^{k\sigma}}\qquad(\text{principal branch; }|xp^{-\sigma}|\le2^{-1/2}).
\]

\begin{reminder}[from Volumes~4, 7, 11 and~14]
\emph{Stone--Weierstrass} (Volume~7, Theorem~4.19). On a compact space, an algebra of continuous complex functions that separates points, contains the constants and is stable under conjugation is dense for the uniform norm.
\emph{The three-series theorem} (Volume~11, Theorem~2.5). A series of independent, centred random variables bounded by $c$ converges almost surely if the sum of their variances is finite.
\emph{$\zeta$ as a short sum} (Volume~14, Lemma~1.5). For integers $k\ge2$, $N\ge1$ and $s=\sigma+it\ne1$, $\sigma>0$: if $N\ge(|s|+k)^{k/(k-1)}$, then $\zeta(s)=\sum_{n\le N}n^{-s}+\frac{N^{1-s}}{s-1}+E$ with $|E|\le\frac{k+1}2N^{-\sigma}$.
\emph{Inversion} (Volume~4, Theorem~6.8). A real random variable whose characteristic function $\varphi$ is integrable has the continuous density $\frac1{2\pi}\int e^{-iux}\varphi(u)\dd u$.
\emph{Riemann--Lebesgue} (Volume~7, Exercise~8.8). $\int f(x)e^{iux}\dd x\to0$ as $|u|\to\infty$ for integrable $f$.
\end{reminder}

\begin{figure}[ht]
\centering\fig{ch06-curve}
\caption{The curve $t\mapsto\zeta(0.7+it)$ for $0\le t\le60$, starting from $\zeta(0.7)=-2.78$ (rust dot). It loops around the origin without passing through it ($\zeta$ has no zeros on this line, if the Riemann hypothesis holds), and by Corollary~\ref{cor:dense} it will eventually pass close to every point of the plane.}
\label{fig:zetacurve}
\end{figure}

\section{The random Euler product}

\begin{proposition}\label{prop:randomeuler}
Almost surely the series $L(\sigma,X)=\sum_p\ell_p(X_p)$ converges, and so does the product $\zeta(\sigma,X)=\prod_p(1-X_pp^{-\sigma})^{-1}=e^{L(\sigma,X)}$, which is therefore non-zero. Moreover $\mathbb E\ell_p(X_p)=0$ and $\mathbb E|\ell_p(X_p)|^2=\sum_{k\ge1}\frac1{k^2p^{2k\sigma}}$.
\end{proposition}
\begin{proof}
Write $\ell_p(X_p)=X_pp^{-\sigma}+r_p$ with $|r_p|\le\sum_{k\ge2}\frac{p^{-k\sigma}}k\le\frac{p^{-2\sigma}}{2(1-2^{-1/2})}\le2p^{-2\sigma}$. Since $2\sigma>1$, $\sum_p|r_p|$ converges. The real parts $\Re(X_pp^{-\sigma})$ are independent, centred, bounded by $1$, with variances $\frac12p^{-2\sigma}$, whose sum is finite; so $\sum_p\Re(X_pp^{-\sigma})$ converges almost surely by the three-series theorem, and likewise the imaginary parts. Hence $L(\sigma,X)$ converges, and $\prod_{p\le y}(1-X_pp^{-\sigma})^{-1}=\exp\sum_{p\le y}\ell_p(X_p)\to e^{L(\sigma,X)}$. Finally $\mathbb EX_p^k=0$ and $\mathbb EX_p^k\overline{X_p^j}=\mathbf 1\{j=k\}$ for $j,k\ge1$, which gives the two moments (the series converges uniformly in $X_p$).
\end{proof}

\begin{lemma}[Kronecker--Weyl for the primes]\label{lem:kw}
Let $p_1,\dots,p_r$ be distinct primes. As $T\to\infty$, $(p_1^{-it},\dots,p_r^{-it})$ converges in law to $(X_{p_1},\dots,X_{p_r})$: for every continuous $F$ on the torus $\mathbb T^r=\{|z_1|=\dots=|z_r|=1\}$, $\mathbb E_TF(p_1^{-it},\dots,p_r^{-it})\to\mathbb EF(X_{p_1},\dots,X_{p_r})$.
\end{lemma}
\begin{proof}
For $m\in\Z^r$ and the monomial $F(z)=z_1^{m_1}\cdots z_r^{m_r}$, $F(p_1^{-it},\dots)=q^{-it}$ with $q=\prod_jp_j^{m_j}$, a positive rational. If $m=0$, both sides are $1$. If $m\ne0$, then $q\ne1$ by unique factorisation, so $|\mathbb E_Tq^{-it}|=|\frac1T\int_T^{2T}e^{-it\log q}\dd t|\le\frac2{T|\log q|}\to0$, while $\mathbb E\prod_jX_{p_j}^{m_j}=0$. By linearity the convergence holds for trigonometric polynomials, which form an algebra that contains the constants, is stable under conjugation ($\overline{z_j}=z_j^{-1}$ on the torus) and separates points. By Stone--Weierstrass they are uniformly dense in $C(\mathbb T^r)$, and a uniform approximation within $\varepsilon$ changes both sides by at most $\varepsilon$.
\end{proof}

\begin{think}
Where did the proof use that the $p_j$ are primes? Try $(2^{-it},4^{-it})$: what is its limit law, and why is it not that of two independent uniform points? (Exercise~\ref{ex:kwfail}.)
\end{think}

\section{The mean square of a mollified zeta}

To pass from finitely many primes to $\zeta$, we show that $\zeta(\sigma+it)$ is close, in mean square, to its Euler product over the primes up to $y$. The tool is a crude mean value theorem for Dirichlet polynomials.

\begin{lemma}\label{lem:crudemvt}
For complex $a_1,\dots,a_N$ and $T>0$,
\[
\Bigl|\int_T^{2T}\Bigl|\sum_{n\le N}a_nn^{-it}\Bigr|^2\dd t-T\sum_{n\le N}|a_n|^2\Bigr|\le4\sum_{m<n\le N}|a_ma_n|\frac n{n-m}.
\]
If moreover $|a_n|\le Kn^{-\sigma}$ with $\frac12<\sigma<1$, the right side is at most $\frac{32K^2}{(1-\sigma)(2-2\sigma)}N^{2-2\sigma}(1+\log N)$.
\end{lemma}
\begin{proof}
Expanding the square, the diagonal terms give $T\sum|a_n|^2$, and for $m<n$ the two terms $(m,n)$ and $(n,m)$ contribute at most $2|a_ma_n||\int_T^{2T}e^{it\log(n/m)}\dd t|\le\frac{4|a_ma_n|}{\log(n/m)}$, and $\log\frac nm\ge1-\frac mn=\frac{n-m}n$. For the second claim, fix $n$ and split the sum over $m<n$. For $m\le\frac n2$, $\frac n{n-m}\le2$ and $\sum_{m\le n/2}m^{-\sigma}\le\frac{n^{1-\sigma}}{1-\sigma}$, which gives at most $\frac{2n^{1-2\sigma}}{1-\sigma}$ after multiplying by $n^{-\sigma}$. For $\frac n2<m<n$, $m^{-\sigma}\le2n^{-\sigma}$ and $\sum_{n/2<m<n}\frac n{n-m}\le n(1+\log n)$, which gives at most $2n^{1-2\sigma}(1+\log n)$. With the factors $K^2$, the sum over $m<n$ is at most $\frac{4K^2n^{1-2\sigma}(1+\log n)}{1-\sigma}$. Finally $\sum_{n\le N}n^{1-2\sigma}\le1+\frac{N^{2-2\sigma}}{2-2\sigma}\le\frac{2N^{2-2\sigma}}{2-2\sigma}$, and $4\cdot4\cdot2=32$.
\end{proof}

Let $A_y(s)=\prod_{p\le y}(1-p^{-s})$, a Dirichlet polynomial: $A_y(s)=\sum_{d\mid P_y}\mu(d)d^{-s}$ with $P_y=\prod_{p\le y}p$.

\begin{proposition}\label{prop:mollified}
For $y\ge2$,
\[
\limsup_{T\to\infty}\ \mathbb E_T\bigl|\zeta(\sigma+it)A_y(\sigma+it)-1\bigr|^2\le\sum_{n>y}\frac1{n^{2\sigma}}.
\]
\end{proposition}
\begin{proof}
Choose $\eta>0$ with $(1+\eta)(2-2\sigma)<1$, which is possible because $\sigma>\frac12$, then an integer $k\ge2$ with $\frac k{k-1}\le1+\eta$, and let $N=\lceil(3T+k)^{1+\eta}\rceil$. For $t\in[T,2T]$ and $T\ge2$, $|s|+k\le3T+k$, so Volume~14, Lemma~1.5 gives $\zeta(s)=Z_N(s)+R(s)$ with $Z_N(s)=\sum_{n\le N}n^{-s}$ and
\[
|R(s)|\le\frac{N^{1-\sigma}}{T}+\frac{k+1}2N^{-\sigma}\longrightarrow0\quad\text{uniformly, since }(1+\eta)(1-\sigma)<1.
\]
Since $|A_y(s)|\le\prod_{p\le y}2=2^{\pi(y)}$, it suffices to show that $\limsup\mathbb E_T|Z_NA_y-1|^2\le\sum_{n>y}n^{-2\sigma}$ (then use $|u+v|^2\le(1+\delta)|u|^2+(1+\frac1\delta)|v|^2$ and let $\delta\to0$).

Multiplying out, $Z_N(s)A_y(s)=\sum_{n\le NP_y}c_nn^{-s}$ with $c_n=\sum\mu(d)$ over the $d\mid\gcd(n,P_y)$ with $\frac nd\le N$. For $n\le N$ the condition $\frac nd\le N$ is automatic, and $c_n=\sum_{d\mid\gcd(n,P_y)}\mu(d)$ is $1$ if $n$ has no prime factor $\le y$ and $0$ otherwise. In particular $c_1=1$, and for $1<n\le N$, $c_n\ne0$ forces $n>y$. For all $n$, $|c_n|\le2^{\pi(y)}$. So $Z_NA_y-1=\sum_{2\le n\le NP_y}a_nn^{-it}$ with $a_n=c_nn^{-\sigma}$, and by Lemma~\ref{lem:crudemvt} with $K=2^{\pi(y)}$,
\[
\mathbb E_T|Z_NA_y-1|^2\le\sum_{y<n\le N}\frac1{n^{2\sigma}}+4^{\pi(y)}\sum_{n>N}\frac1{n^{2\sigma}}+\frac{C_{\sigma}4^{\pi(y)}}T(NP_y)^{2-2\sigma}(1+\log NP_y).
\]
As $T\to\infty$ with $y$ fixed, $N\to\infty$, so the middle term tends to $0$, and the last is $O_{y,\sigma}(T^{(1+\eta)(2-2\sigma)-1}\log T)\to0$ by the choice of $\eta$.
\end{proof}

The proof shows the mechanism clearly: $\zeta A_y-1$ is a Dirichlet series over the integers with no prime factor up to $y$, and its mean square is the sum of the squares of its coefficients, as if the $n^{-it}$ were orthogonal. The only role of $\sigma>\frac12$ is to make that sum, $\sum_{n>y}n^{-2\sigma}$, finite and small.

\section{The theorem of Bohr and Jessen}

\begin{theorem}[Bohr and Jessen, 1930--1932]\label{thm:bohrjessen}
Let $\frac12<\sigma<1$. As $T\to\infty$, $\zeta(\sigma+it)$ with $t$ uniform on $[T,2T]$ converges in law to $\zeta(\sigma,X)$: for every bounded Lipschitz $g:\C\to\R$,
\[
\mathbb E_Tg\bigl(\zeta(\sigma+it)\bigr)\longrightarrow\mathbb Eg\bigl(\zeta(\sigma,X)\bigr).
\]
Consequently $\log|\zeta(\sigma+it)|$ converges in law to $\Re L(\sigma,X)=\log|\zeta(\sigma,X)|$.
\end{theorem}
\begin{proof}
Let $g$ be bounded by $1$ and $1$-Lipschitz (the general case follows by scaling). For $y\ge2$ let $E_y(t)=A_y(\sigma+it)^{-1}=\prod_{p\le y}(1-p^{-\sigma-it})^{-1}$ and $\zeta_y(X)=\prod_{p\le y}(1-X_pp^{-\sigma})^{-1}$.

\emph{Finitely many primes.} $E_y(t)=\Phi(p^{-it})_{p\le y}$ for the continuous function $\Phi(z)=\prod_{p\le y}(1-z_pp^{-\sigma})^{-1}$ on the torus, so by Lemma~\ref{lem:kw}, $\mathbb E_Tg(E_y)\to\mathbb Eg(\zeta_y(X))$ and $\mathbb E_T|E_y|^2\to\mathbb E|\zeta_y(X)|^2$. By independence and Parseval, $\mathbb E|1-X_pp^{-\sigma}|^{-2}=\mathbb E|\sum_kX_p^kp^{-k\sigma}|^2=\sum_kp^{-2k\sigma}=(1-p^{-2\sigma})^{-1}$, so $\mathbb E|\zeta_y(X)|^2=\prod_{p\le y}(1-p^{-2\sigma})^{-1}\le\zeta(2\sigma)$.

\emph{Comparison.} Since $\zeta=E_y\cdot(\zeta A_y)$, by Cauchy--Schwarz,
\[
\bigl|\mathbb E_Tg(\zeta)-\mathbb E_Tg(E_y)\bigr|\le\mathbb E_T|\zeta-E_y|=\mathbb E_T\bigl[|E_y||\zeta A_y-1|\bigr]\le\bigl(\mathbb E_T|E_y|^2\bigr)^{1/2}\bigl(\mathbb E_T|\zeta A_y-1|^2\bigr)^{1/2}.
\]
By the first step and Proposition~\ref{prop:mollified},
\[
\limsup_{T\to\infty}\bigl|\mathbb E_Tg(\zeta)-\mathbb Eg(\zeta_y(X))\bigr|\le\zeta(2\sigma)^{1/2}\varepsilon_y,\qquad\varepsilon_y=\Bigl(\sum_{n>y}n^{-2\sigma}\Bigr)^{1/2}.
\]

\emph{All primes.} By Proposition~\ref{prop:randomeuler}, $\zeta_y(X)\to\zeta(\sigma,X)$ almost surely, so $\mathbb Eg(\zeta_y(X))\to\mathbb Eg(\zeta(\sigma,X))$ by dominated convergence. Letting $y\to\infty$, $\limsup_T|\mathbb E_Tg(\zeta)-\mathbb Eg(\zeta(\sigma,X))|\le0$.

\emph{Logarithms.} The function $z\mapsto\log|z|$ is continuous on $\C\setminus\{0\}$ and $\zeta(\sigma,X)\ne0$ almost surely. For bounded Lipschitz $h:\R\to\R$ and $\delta>0$, let $g_\delta(z)=h(\log\max(|z|,\delta))$, which is bounded and Lipschitz on $\C$. Then $|\mathbb E_Th(\log|\zeta|)-\mathbb E_Tg_\delta(\zeta)|\le2\|h\|_\infty P_T(|\zeta|<\delta)$ and similarly for $\zeta(\sigma,X)$. Since $P_T(|\zeta|<\delta)\le\mathbb E_T\phi_\delta(\zeta)\to\mathbb E\phi_\delta(\zeta(\sigma,X))\le P(|\zeta(\sigma,X)|<2\delta)$ for the Lipschitz function $\phi_\delta(z)=\min(1,(2-\frac{|z|}\delta)^+)$, and $P(|\zeta(\sigma,X)|<2\delta)\to0$ as $\delta\to0$, the claim follows.
\end{proof}

A remarkable feature of the proof is that it never mentions the zeros of $\zeta$: if $\zeta$ had zeros with real part $\sigma$, they would occupy a set of $t$ of measure tending to $0$ in proportion (Volume~14, Corollary~1.8 already says they are rare), and the theorem would not see them. Figure~\ref{fig:bohrjessen} compares the two laws numerically.

\begin{figure}[ht]
\centering\fig{ch06-density}
\caption{Histogram of $\log|\zeta(0.7+it)|$ at $10\,000$ random points $t\in[10^5,2\cdot10^5]$ (bars), against the density of $\log|\zeta(0.7,X)|$ estimated from $40\,000$ samples of the random Euler product over $p\le10^5$ (rust). The Kolmogorov--Smirnov distance between the two samples is $0.0095$, within sampling error; the variances are $0.525$ and $0.531$, and the exact variance $\frac12\sum_p\sum_k\frac1{k^2p^{1.4k}}$ is $0.529$ (Exercise~\ref{ex:varbj}).}
\label{fig:bohrjessen}
\end{figure}

\section{Every value, and a smooth density}

\begin{lemma}\label{lem:support}
For every $w\in\C$ and $\varepsilon>0$, $P(|L(\sigma,X)-w|<\varepsilon)>0$.
\end{lemma}
\begin{proof}
Write $\ell_p(x)=xp^{-\sigma}+r_p(x)$ with $|r_p(x)|\le2p^{-2\sigma}$, as in Proposition~\ref{prop:randomeuler}. Choose $y_0$ such that, for $p>y_0$, $p^{-\sigma}<\frac\varepsilon4$, and such that $\sum_{p>y_0}2p^{-2\sigma}<\frac\varepsilon4$ and $\sum_{p>y_0}\mathbb E|\ell_p(X_p)|^2<\frac{\varepsilon^2}{32}$.

\emph{A deterministic choice.} Let $x_p=1$ for $p\le y_0$, and $w'=w-\sum_{p\le y_0}\ell_p(1)$. For the primes $p>y_0$ in increasing order, let $D$ be the current value of $w'-\sum_{y_0<q<p}\ell_q(x_q)$ and choose $x_p=\frac D{|D|}$ (or $x_p=1$ if $D=0$). If $|D|\ge p^{-\sigma}$, the new value has modulus at most $|D|-p^{-\sigma}+2p^{-2\sigma}$. Since $\sum_pp^{-\sigma}=\infty$ for $\sigma<1$ and $\sum_pp^{-2\sigma}<\infty$, this cannot happen for every $p$; at the first prime $p=y_1$ with $|D|<p^{-\sigma}$, the new value has modulus at most $p^{-\sigma}+2p^{-2\sigma}<\frac\varepsilon2$. So $|\sum_{p\le y_1}\ell_p(x_p)-w|<\frac\varepsilon2$.

\emph{Positive probability.} The map $(z_p)_{p\le y_1}\mapsto\sum_{p\le y_1}\ell_p(z_p)$ is continuous on the torus, so there is $\delta>0$ such that $|\sum_{p\le y_1}\ell_p(X_p)-w|<\frac{3\varepsilon}4$ whenever $|X_p-x_p|<\delta$ for all $p\le y_1$. This event $B$ has probability $\prod_{p\le y_1}P(|X_p-x_p|<\delta)>0$. The tail $R=\sum_{p>y_1}\ell_p(X_p)$ is independent of $B$, centred, with $\mathbb E|R|^2=\sum_{p>y_1}\mathbb E|\ell_p(X_p)|^2<\frac{\varepsilon^2}{32}$ (independent centred terms), so $P(|R|\ge\frac\varepsilon4)\le\frac12$ by Chebyshev. Hence $P(|L(\sigma,X)-w|<\varepsilon)\ge\frac12P(B)>0$.
\end{proof}

\begin{corollary}[Bohr and Courant, 1914]\label{cor:dense}
Let $\frac12<\sigma<1$. For every $z\in\C$ and $\varepsilon>0$ there are arbitrarily large $t$ with $|\zeta(\sigma+it)-z|<\varepsilon$. In fact the set of such $t$ in $[T,2T]$ has measure at least $cT$ for some $c=c(z,\varepsilon,\sigma)>0$ and all large $T$.
\end{corollary}
\begin{proof}
We may assume $z\ne0$ (approach $0$ through non-zero points) and $\varepsilon<|z|$. Write $z=e^w$. By continuity of $\exp$ at $w$ there is $\varepsilon'>0$ with $|e^v-z|<\frac\varepsilon2$ when $|v-w|<\varepsilon'$, so $q=P(|\zeta(\sigma,X)-z|<\frac\varepsilon2)>0$ by Lemma~\ref{lem:support}. Let $g(u)=\min(1,(2-\frac{2|u-z|}\varepsilon)^+)$, which is Lipschitz, equals $1$ on $|u-z|\le\frac\varepsilon2$ and vanishes off $|u-z|<\varepsilon$. By Theorem~\ref{thm:bohrjessen}, $P_T(|\zeta(\sigma+it)-z|<\varepsilon)\ge\mathbb E_Tg(\zeta)\to\mathbb Eg(\zeta(\sigma,X))\ge q$, so for large $T$ the measure of the set is at least $\frac q2T$.
\end{proof}

\begin{beauty}
Bohr's theorem is about a single deterministic function, and its proof is a statement about probability: the values of $\zeta$ come close to $z$ because a random Euler product does so with positive probability, and the primes, through Kronecker--Weyl, are random enough. No explicit $t$ is ever found. In 1975 Sergei Voronin went much further: every non-vanishing holomorphic function on a small disc inside the strip is approximated by a vertical translate $\zeta(s+it)$ of zeta. Bhaskar Bagchi showed in 1981 that this \emph{universality} is again a support theorem for a random Euler product, now viewed as a random holomorphic function.
\end{beauty}

\begin{theorem}[Jessen and Wintner, 1935]\label{thm:smoothdensity}
The random variable $\log|\zeta(\sigma,X)|$ has a density of class $C^\infty$ on $\R$. Its characteristic function $\varphi(u)$ satisfies $|\varphi(u)|\le\exp(-c|u|^{1/\sigma}/\log|u|)$ for $|u|\ge u_0$, with $c,u_0>0$ depending on $\sigma$.
\end{theorem}
\begin{proof}
\emph{Factorisation.} By independence and the almost sure convergence of $\sum_p\Re\ell_p(X_p)$, $\varphi(u)=\mathbb Ee^{iu\log|\zeta(\sigma,X)|}=\prod_p\varphi_p(u)$ with $\varphi_p(u)=\mathbb Ee^{iu\Re\ell_p(X_p)}$, and $|\varphi_p(u)|\le1$ (dominated convergence for the partial products).

\emph{One prime.} Let $J(v)=\mathbb Ee^{iv\Re X_p}=\frac1{2\pi}\int_0^{2\pi}e^{iv\cos\theta}\dd\theta=\int_{-1}^1e^{ivx}\frac{\dd x}{\pi\sqrt{1-x^2}}$ (the real part of a uniform point of the circle has the arcsine density on $[-1,1]$). By Riemann--Lebesgue, $J(v)\to0$; fix $A\ge1$ with $|J(v)|\le\frac14$ for $v\ge A$. Since $|e^{ia}-e^{ib}|\le|a-b|$ and $|\Re\ell_p(X_p)-\Re X_pp^{-\sigma}|\le2p^{-2\sigma}$,
\[
|\varphi_p(u)|\le|J(|u|p^{-\sigma})|+2|u|p^{-2\sigma}.
\]

\emph{Many primes.} Let $u\ge u_0$ (the case $u<0$ follows from $\varphi(-u)=\overline{\varphi(u)}$), and consider the primes with $A\le up^{-\sigma}\le2A$, that is $(\frac u{2A})^{1/\sigma}\le p\le(\frac uA)^{1/\sigma}$. For these, $2up^{-2\sigma}\le4Ap^{-\sigma}\le\frac14$ once $u$ is large, so $|\varphi_p(u)|\le\frac12$. By the prime number theorem their number is at least $c_1\frac{u^{1/\sigma}}{\log u}$ for $u\ge u_0$, with $c_1=c_1(A,\sigma)>0$, because the interval is $[aU,bU]$ with $U=u^{1/\sigma}$ and fixed $0<a<b$. Bounding the other factors by $1$, $|\varphi(u)|\le2^{-c_1u^{1/\sigma}/\log u}$.

\emph{Smoothness.} Hence $\int|u|^k|\varphi(u)|\dd u<\infty$ for every $k$. By Volume~4, Theorem~6.8, the law has the continuous density $f(v)=\frac1{2\pi}\int e^{-iuv}\varphi(u)\dd u$, and differentiating $k$ times under the integral sign (justified by dominated convergence with the integrable bound $|u|^k|\varphi(u)|$) shows that $f$ is $C^k$ for every $k$.
\end{proof}

\begin{remark}
These are quoted for culture; see the Sources. (a) Jessen and Wintner (1935) proved that the law of the complex variable $L(\sigma,X)$ also has a smooth density in the plane. (b) The upper tail is much thinner than Gaussian: Lamzouri proved that the proportion of $t\in[T,2T]$ with $\log|\zeta(\sigma+it)|\ge\tau$ is $\exp\bigl(-A(\sigma)\tau^{1/(1-\sigma)}(\log\tau)^{\sigma/(1-\sigma)}(1+o(1))\bigr)$, uniformly for $\tau$ up to a constant times $(\log T)^{1-\sigma}/\log\log T$, the same tail as for the random Euler product. Montgomery (1977) had shown that values this large do occur. (c) Voronin's universality theorem (1975) and Bagchi's thesis (1981) described in the box above.
\end{remark}


\section{Selberg's central limit theorem}\label{ch:selbergclt}

\begin{center}\itshape
On the critical line the Euler product diverges, the zeros sit in the way,\\
and $\log|\zeta|$ plunges to $-\infty$ infinitely often.\\
And yet, seen at a random height, it is a Gaussian with an explicit variance.
\end{center}

In April 1940 the war came to Norway, and a twenty-two-year-old student called Atle Selberg (1917--2007) stopped doing mathematics. He fought with the Norwegian forces against the invasion and was taken prisoner at the camp of Trandum. Released, he defended his doctorate, was arrested by the Germans again immediately after the examination, and, once out of prison, spent the rest of the war in Gj\o{}vik, working on the Riemann zeta function on his own. Two results of those isolated years would have been enough for a career. In 1942 he proved that a positive proportion of the zeros of $\zeta$ lie on the critical line, where Hardy had only found infinitely many. And in a long paper of 1946, published in a Norwegian journal that few people read, he proved that $\log|\zeta(\frac12+it)|$, for $t$ chosen at random, is approximately Gaussian with variance $\frac12\log\log t$. Four years later he received a Fields Medal, for other work. The central limit theorem waited until a survey talk of 1989 to become widely known.

Section~\ref{ch:bohrjessen} explains why the theorem should be true. For $\frac12<\sigma<1$ fixed, $\log|\zeta(\sigma+it)|$ has a limit law, the law of a random Euler product, and its variance $\frac12\sum_p\sum_k\frac1{k^2p^{2k\sigma}}$ is $\frac12\log\frac1{\sigma-\frac12}+O(1)$ (Exercise~\ref{ex:varbj}). The function $\log\zeta$ at height $T$ only ``sees'' the primes up to about $T$, and the Euler product over those primes behaves as if $\sigma-\frac12$ were replaced by $\frac1{\log T}$. Substitute, and the variance becomes $\frac12\log\log T$. On the critical line itself the limit law disappears (the variance tends to infinity), and what survives is its shape: the sum of many independent small contributions, one for each prime, is Gaussian.

Turning that heuristic into a proof is hard for a reason that Section~\ref{ch:bohrjessen} never met: on $\Re s=\frac12$, $\zeta$ has zeros, infinitely many of them, and at each one $\log|\zeta|=-\infty$. Selberg controlled them with a zero density estimate of his own. This chapter follows instead a proof found in 2015 by Maksym Radziwi\l\l{} and Kannan Soundararajan, which needs almost nothing about the zeros: one integral, and the count $N(T)$. Its strategy has four steps, each a section:
\begin{enumerate}
\item move from $\frac12$ to $\sigma_0=\frac12+\frac W{\log T}$, slightly to the right, at a cost controlled by the zeros on average;
\item at $\sigma_0$, the prime sum $P(\sigma_0+it)=\sum_{n\le X}\frac{\Lambda(n)}{n^{\sigma_0+it}\log n}$ is Gaussian, by the method of moments;
\item a short Dirichlet polynomial $M$, the \emph{mollifier}, is close to $e^{-P}$;
\item $\zeta M$ is close to $1$ in mean square, by a twisted second moment of $\zeta$.
\end{enumerate}
Then $\log|\zeta(\frac12+it)|\approx\log|\zeta(\sigma_0+it)|\approx-\log|M(\sigma_0+it)|\approx\operatorname{Re}P(\sigma_0+it)$, which is Gaussian.

\begin{reminder}[from Volumes~4, 5, 6 and~14]
\emph{Zeros} (Volume~6, Proposition~6.2, Theorems~6.7 and~6.9, Corollary~6.10). $\xi(s)=G(s)\zeta(s)$ with $G(s)=\frac12s(s-1)\pi^{-s/2}\Gamma(\frac s2)$ is entire, $\xi(s)=\xi(1-s)$, $\xi(\overline s)=\overline{\xi(s)}$, and $|\xi(s)|\le e^{C|s|\log|s|}$ for $|s|\ge2$. Its zeros $\rho=\beta+i\gamma$ are the non-trivial zeros of $\zeta$, $0<\beta<1$, and $\rho\mapsto1-\overline\rho$ permutes them (with multiplicity). Hadamard: $\frac{\xi'}\xi(s)=B+\sum_\rho(\frac1{s-\rho}+\frac1\rho)$, absolutely convergent, and $\sum_\rho|\rho|^{-2}<\infty$. $N(T)=\frac T{2\pi}\log\frac T{2\pi e}+\frac78+O(\log T)$, there are $O(\log T)$ zeros with $T<\gamma\le T+1$, and $\sum_\rho\frac1{1+(t-\gamma)^2}=O(\log|t|)$ for $|t|\ge2$.
\emph{Stirling} (Volume~6, Theorem~6.1). For $|\arg w|\le\pi-\delta$, $|w|\ge1$: $\log\Gamma(w)=(w-\frac12)\log w-w+\frac12\log2\pi+O_\delta(\frac1{|w|})$ and $\frac{\Gamma'}\Gamma(w)=\log w+O_\delta(\frac1{|w|})$.
\emph{$\zeta$ to the right of $0$} (Volume~5, Proposition~12.6, and Volume~6, Exercise~6.6). $\zeta(s)=\frac1{s-1}+\gamma+O(s-1)$ near $1$, and for $\Re s>0$, $\zeta(s)=\frac1{s-1}+1-s\int_1^\infty\{u\}u^{-s-1}\dd u$, so $|\zeta(s)-\frac1{s-1}|\le1+\frac{|s|}{\Re s}$.
\emph{$\zeta$ as a short sum} (Volume~14, Lemma~1.5), as recalled in Section~\ref{ch:bohrjessen}.
\emph{Mertens} (Volume~6, Theorems~1.3 and~1.5). $\sum_{p\le x}\frac{\log p}p=\log x+O(1)$ and $\sum_{p\le x}\frac1p=\log\log x+M+O(\frac1{\log x})$.
\emph{Moments} (Volume~6, Lemma~9.6). If $\mathbb E[Y_N^k]\to\mathbb E[Z^k]$ for every $k\ge1$, with $Z$ standard Gaussian, then $Y_N\to Z$ in law. \emph{Slutsky} (Volume~4, Exercise~6.1): if $Y_N\to Z$ in law and $Y_N'-Y_N\to0$ in probability, then $Y_N'\to Z$ in law.
\end{reminder}

\begin{figure}[ht]
\centering\fig{ch07-hist}
\caption{Histogram of $\log|\zeta(\frac12+it)|$ at $8000$ random points $t\in[10^{10},2\cdot10^{10}]$, computed with the Riemann--Siegel formula. Mean $0.003$, variance $2.20$. Selberg's variance $\frac12\log\log T=1.57$ (dashed) is far too small at this height; the random matrix prediction $\frac12(\log N+\gamma+1)=2.32$ of Keating and Snaith, with $N=\log\frac t{2\pi}$ at the middle of the range, is much closer (rust). The left tail is long: the smallest value in the sample is $-10.8$, at a point very close to a zero.}
\label{fig:selberghist}
\end{figure}

Throughout the chapter, $t$ is uniform on $[T,2T]$ and $P_T$, $\mathbb E_T$ refer to this choice; $\operatorname{meas}$ is Lebesgue measure. The theorem to be proved is the following.

\begin{theorem}[Selberg, 1946]\label{thm:selberg}
For every real $a$,
\[
P_T\Bigl(\log|\zeta(\tfrac12+it)|\le a\sqrt{\tfrac12\log\log T}\Bigr)\longrightarrow\Phi(a)=\frac1{\sqrt{2\pi}}\int_{-\infty}^ae^{-u^2/2}\dd u\qquad(T\to\infty).
\]
\end{theorem}

The zeros of $\zeta$ on the line form a countable set, so $\log|\zeta(\frac12+it)|$ is defined for almost every $t$. Figure~\ref{fig:selberghist} shows the law at height $10^{10}$, and a warning: the convergence is slow, with errors of relative size $1/\log\log T$, so that the theorem says little about the numbers we can compute. We return to this after the proof.

\section{Off the line, at a cost}

The first step is a short computation with the zeros. For $\rho=\beta+i\gamma$ and $\sigma>\frac12$, put
\[
f_\rho(t)=\log\Bigl|\frac{\sigma+it-\rho}{\frac12+it-\rho}\Bigr|=\frac12\log\frac{(\sigma-\beta)^2+(t-\gamma)^2}{(\frac12-\beta)^2+(t-\gamma)^2}.
\]

\begin{lemma}\label{lem:logint}
For real $a,b$, $\displaystyle\int_{-\infty}^\infty\frac12\log\frac{a^2+u^2}{b^2+u^2}\dd u=\pi(|a|-|b|)$. In particular $\int_\R|f_\rho(t)|\dd t\le\pi(\sigma-\frac12)$ for every $\rho$ with $0<\beta<1$.
\end{lemma}
\begin{proof}
The integrand is even in $a$ and in $b$, so take $a,b\ge0$. It is integrable: near $u=0$ it has at most a logarithmic singularity, and for large $u$ it is $O(\frac{|a^2-b^2|}{u^2})$. Let $F(a)$ be the integral, for fixed $b$. For $a>0$, differentiation under the integral sign (dominated by $\frac{2a}{a^2+u^2}$ locally uniformly in $a$) gives $F'(a)=\int\frac a{a^2+u^2}\dd u=\pi$. As $F(b)=0$ and $F$ is continuous on $[0,\infty[$ (dominated convergence), $F(a)=\pi(a-b)$. For $f_\rho$, the sign of the integrand is that of $|\sigma-\beta|-|\frac12-\beta|$ for every $t$, so $\int|f_\rho|=\pi\bigl||\sigma-\beta|-|\frac12-\beta|\bigr|\le\pi|\sigma-\frac12|$ by the triangle inequality.
\end{proof}

\begin{proposition}\label{prop:offline}
There are absolute constants $C$ and $T_0$ such that for $T\ge T_0$ and $\frac12<\sigma\le1$,
\[
\int_T^{2T}\bigl|\log|\zeta(\tfrac12+it)|-\log|\zeta(\sigma+it)|\bigr|\dd t\le C\Bigl(\sigma-\frac12\Bigr)T\log T.
\]
\end{proposition}
\begin{proof}
\emph{The zeros cancel in pairs at $\frac12$.} For $s=u+it$ not a zero, $\operatorname{Re}\frac{\xi'}\xi(s)=\operatorname{Re}B+\sum_\rho\operatorname{Re}\frac1\rho+\sum_\rho\frac{u-\beta}{(u-\beta)^2+(t-\gamma)^2}$. Both sums converge absolutely: $\operatorname{Re}\frac1\rho=\frac\beta{|\rho|^2}\le\frac1{|\rho|^2}$, and the terms of the second are $O(\frac1{1+(t-\gamma)^2})$ except for the finitely many $\rho$ with $|t-\gamma|<1$. Let $t$ be such that $\xi(\frac12+it)\ne0$. The function $v\mapsto|\xi(\frac12+v+it)|$ is even, because $\xi(\frac12-v+it)=\xi(\frac12+v-it)=\overline{\xi(\frac12+v+it)}$; so its logarithm has derivative $0$ at $v=0$, that is $\operatorname{Re}\frac{\xi'}\xi(\frac12+it)=0$. And $\sum_\rho\frac{\frac12-\beta}{(\frac12-\beta)^2+(t-\gamma)^2}=0$, since $\rho\mapsto1-\overline\rho$ changes $\beta$ into $1-\beta$, the term into its opposite, and keeps $\gamma$. Hence $\operatorname{Re}B+\sum_\rho\operatorname{Re}\frac1\rho=0$, and for every $t$ that is not the ordinate of a zero, integrating in $u$ from $\frac12$ to $\sigma$ (the series converges uniformly in $u$, because $t\ne\gamma$ for every $\rho$),
\[
\log|\xi(\sigma+it)|-\log|\xi(\tfrac12+it)|=\sum_\rho\int_{1/2}^\sigma\frac{u-\beta}{(u-\beta)^2+(t-\gamma)^2}\dd u=\sum_\rho f_\rho(t).
\]

\emph{The Gamma factor.} $\frac{G'}G(s)=\frac1s+\frac1{s-1}-\frac12\log\pi+\frac12\frac{\Gamma'}\Gamma(\frac s2)$, and by Stirling, for $T\le t\le2T$ and $\frac12\le u\le1$, $\operatorname{Re}\frac{G'}G(u+it)=\frac12\log\frac t{2\pi}+O(\frac1t)$. So $\bigl|\log|G(\sigma+it)|-\log|G(\frac12+it)|\bigr|\le(\sigma-\frac12)\log T$ for $T$ large. Since $\log|\zeta|=\log|\xi|-\log|G|$, it remains to bound $\int_T^{2T}\sum_\rho|f_\rho|$.

\emph{Counting zeros.} Let $d_\rho$ be the distance from $\gamma$ to $[T,2T]$. If $d_\rho\ge1$, then for $t\in[T,2T]$, writing $D=(\sigma-\beta)^2-(\frac12-\beta)^2=(\sigma-\frac12)(\sigma+\frac12-2\beta)$, so that $|D|\le\frac32(\sigma-\frac12)$, and using $|\log(1+x)|\le\frac{|x|}{\min(1,1+x)}$ for $x>-1$,
\[
|f_\rho(t)|=\frac12\Bigl|\log\Bigl(1+\frac D{(\frac12-\beta)^2+(t-\gamma)^2}\Bigr)\Bigr|\le\frac{|D|}{2(t-\gamma)^2}\le\frac{\sigma-\frac12}{d_\rho^2}.
\]
With Lemma~\ref{lem:logint}, $\int_T^{2T}|f_\rho|\le(\sigma-\frac12)\min(\pi,\frac T{d_\rho^2})$ for every $\rho$ (with $\frac T{0}=\infty$). There are $N(2T+1)-N(T-1)=O(T\log T)$ zeros with $d_\rho<1$ and $\gamma>0$, and for $k\ge1$, $O(\log(T+k))$ zeros with $k\le d_\rho<k+1$ and $\gamma>0$; the zeros with $\gamma<0$ are the conjugates of these, with $d_\rho\ge T$. Altogether
\[
\int_T^{2T}\sum_\rho|f_\rho|\le C\Bigl(\sigma-\frac12\Bigr)\Bigl(T\log T+\sum_{k\ge1}\log(T+k)\min\Bigl(\pi,\frac T{k^2}\Bigr)\Bigr)\le C'\Bigl(\sigma-\frac12\Bigr)T\log T,
\]
the sum over $k$ being $O(\sqrt T\log T)$ (split at $k=\sqrt T$). The interchange of $\sum$ and $\int$ is Tonelli's theorem.
\end{proof}

From now on we fix the parameters of Radziwi\l\l{} and Soundararajan. Write $L_2=\log\log T$, $L_3=\log\log\log T$, and
\[
W=L_3^4,\qquad\sigma_0=\frac12+\frac W{\log T},\qq X=T^{1/L_3^2},\qquad Y=T^{1/L_2^2}.
\]
So $\sigma_0-\frac12$ is a little more than $\frac1{\log T}$, and $Y<X$ are both $T^{o(1)}$, with $X$ much larger than $Y$.

\begin{corollary}\label{cor:offline}
For every $\varepsilon>0$, $P_T\bigl(|\log|\zeta(\frac12+it)|-\log|\zeta(\sigma_0+it)||>\varepsilon\sqrt{L_2}\bigr)\to0$.
\end{corollary}
\begin{proof}
By Markov's inequality and Proposition~\ref{prop:offline}, the probability is at most $\frac{C(\sigma_0-\frac12)\log T}{\varepsilon\sqrt{L_2}}=\frac{CL_3^4}{\varepsilon\sqrt{L_2}}\to0$.
\end{proof}

\begin{think}
The corollary needs $W=o(\sqrt{L_2})$. Section~4 will need $(\sigma_0-\frac12)\log X\to\infty$, that is $W\gg\frac{\log T}{\log X}=L_3^2$. Why does each condition appear? Which of them is about the zeros, and which about the primes?
\end{think}

\section{The prime sum is Gaussian}

Two Dirichlet polynomials $A(t)=\sum a_nn^{-it}$ and $B(t)=\sum b_nn^{-it}$ (finite sums) are almost orthogonal over a long interval when their lengths are small: this is the engine of the whole chapter.

\begin{lemma}\label{lem:mvt2}
For finitely supported $(a_n)$, $(b_n)$ and $T>0$,
\[
\Bigl|\int_T^{2T}A(t)\overline{B(t)}\dd t-T\sum_na_n\overline{b_n}\Bigr|\le2\Bigl(\sum_n|a_n|n\Bigr)\Bigl(\sum_n|b_n|n\Bigr).
\]
\end{lemma}
\begin{proof}
Expand. The terms $m=n$ give $T\sum a_n\overline{b_n}$. For $m<n$, $\bigl|\int_T^{2T}(n/m)^{it}\dd t\bigr|\le\frac2{\log(n/m)}$ and $\log\frac nm\ge1-\frac mn\ge\frac1n$, so the term $(m,n)$ or $(n,m)$ is at most $2|a_mb_n|\max(m,n)\le2|a_m|m|b_n|n$.
\end{proof}

Let
\[
P(s)=\sum_{2\le n\le X}\frac{\Lambda(n)}{n^s\log n},\qquad P_0(s)=\sum_{p\le X}\frac1{p^s},\qquad v_T=\sum_{p\le X}\frac1{p^{2\sigma_0}}.
\]
$P$ is the truncation at $X$ of $\log\zeta(s)=\sum_n\frac{\Lambda(n)}{n^s\log n}$ ($\Re s>1$), and $P_0$ keeps only the primes. The following elementary lemma measures sums over primes with a damping factor $p^{-\delta}$; it will serve several times.

\begin{lemma}\label{lem:dampedmertens}
Let $E_1(x)=\int_x^\infty e^{-v}\frac{\dd v}v$, so that $E_1(x)\le1+\log\frac1x$ for $0<x\le1$ and $E_1(x)\le\frac{e^{-x}}x$ for $x\ge1$. There is an absolute $C$ such that for $y\ge2$ and $0<\delta\le1$,
\[
\sum_{p>y}\frac1{p^{1+\delta}}\le E_1(\delta\log y)+\frac C{\log y}.
\]
\end{lemma}
\begin{proof}
The bounds on $E_1$: for $x\le1$, $E_1(x)\le\int_x^1\frac{\dd v}v+\int_1^\infty e^{-v}\dd v\le\log\frac1x+1$; for $x\ge1$, $E_1(x)\le\frac1x\int_x^\infty e^{-v}\dd v$. Write $S(u)=\sum_{p\le u}\frac1p=\log\log u+M+R(u)$ with $|R(u)|\le\frac{C_1}{\log u}$ (Mertens). By partial summation, for $Z>y$,
\[
\sum_{y<p\le Z}\frac{p^{-\delta}}p=\int_y^Z\frac{u^{-\delta}\dd u}{u\log u}+\bigl[u^{-\delta}R(u)\bigr]_y^Z+\delta\int_y^Zu^{-\delta-1}R(u)\dd u.
\]
The first integral is at most $\int_{\delta\log y}^\infty e^{-v}\frac{\dd v}v=E_1(\delta\log y)$ (put $u=e^{v/\delta}$), the bracket is at most $\frac{2C_1}{\log y}$, and the last term at most $\frac{C_1}{\log y}\delta\int_y^\infty u^{-\delta-1}\dd u\le\frac{C_1}{\log y}$. Let $Z\to\infty$.
\end{proof}

\begin{proposition}\label{prop:primesgauss}
$v_T=L_2+O(\log W)$, and as $T\to\infty$,
\[
\frac{\operatorname{Re}P(\sigma_0+it)}{\sqrt{v_T/2}}\longrightarrow\mathcal N(0,1)\quad\text{in law}.
\]
\end{proposition}
\begin{proof}
\emph{The variance.} Let $\delta=2\sigma_0-1=\frac{2W}{\log T}$. On one side $v_T\le\sum_{p\le X}\frac1p=\log\log X+O(1)\le L_2+O(1)$. On the other, with $Z=T^{1/W}\le X$ and $p^{-\delta}=e^{-\delta\log p}\ge1-\delta\log p$,
\[
v_T\ge\sum_{p\le Z}\frac1p-\delta\sum_{p\le Z}\frac{\log p}p=\log\log Z-\delta\log Z+O(1)=L_2-\log W+O(1).
\]

\emph{Moments of $P_0$.} Write $s=\sigma_0+it$ and $P_0(s)^j=\sum_na_j(n)n^{-s}$, where $a_j(n)$ is the number of ordered $j$-tuples of primes $\le X$ with product $n$. Then $\sum_na_j(n)n^{1-\sigma_0}\le X^{j/2}\pi(X)^j\le X^{2j}$, and Lemma~\ref{lem:mvt2} gives, for fixed $j,\ell\ge0$,
\[
\frac1T\int_T^{2T}P_0(s)^j\overline{P_0(s)}^\ell\dd t=\sum_n\frac{a_j(n)a_\ell(n)}{n^{2\sigma_0}}+O\Bigl(\frac{X^{2(j+\ell)}}T\Bigr).
\]
The error tends to $0$ since $X^{2(j+\ell)}=T^{2(j+\ell)/L_3^2}$. If $j\ne\ell$ the sum vanishes, because $a_j(n)\ne0$ forces $\Omega(n)=j$. If $j=\ell$, it equals $\sum a_j(p_1\cdots p_j)(p_1\cdots p_j)^{-2\sigma_0}$ over all $j$-tuples. For tuples of distinct primes $a_j=j!$, and $a_j\le j!$ always; the tuples with a repeated prime have total weight at most $\binom j2v_T^{j-2}\sum_pp^{-4\sigma_0}\le j^2v_T^{j-2}$. So the sum is $j!\,v_T^j+O_j(v_T^{j-2})$.

Now $(\operatorname{Re}P_0)^k=2^{-k}\sum_{j=0}^k\binom kjP_0^j\overline{P_0}^{k-j}$, and only $j=k-j$ survives:
\[
\mathbb E_T\bigl[(\operatorname{Re}P_0(s))^k\bigr]=\begin{cases}2^{-k}\binom k{k/2}(k/2)!\,v_T^{k/2}+O_k(v_T^{k/2-2})+o(1),&k\text{ even},\\o(1),&k\text{ odd}.\end{cases}
\]
Since $2^{-k}\binom k{k/2}(k/2)!=2^{-k/2}(k-1)!!$ (Exercise~\ref{ex:half}) and $v_T\to\infty$, the moments of $\operatorname{Re}P_0(s)/\sqrt{v_T/2}$ tend to those of $\mathcal N(0,1)$, and Volume~6, Lemma~9.6, gives convergence in law.

\emph{From $P_0$ to $P$.} $P-P_0=Q+R$ with $Q(s)=\frac12\sum_{p\le\sqrt X}p^{-2s}$ and $|R(s)|\le\sum_p\sum_{k\ge3}\frac{p^{-k/2}}k\le\sum_p\frac{p^{-3/2}}{1-p^{-1/2}}=O(1)$. By Lemma~\ref{lem:mvt2}, $\int_T^{2T}|Q|^2\le\frac T4\sum_pp^{-4\sigma_0}+2X^3\le T$ for $T$ large, so $P_T(|Q|>A)\le A^{-2}$. Hence $(P-P_0)/\sqrt{v_T}\to0$ in probability, and Slutsky's lemma concludes.
\end{proof}

The factor $\frac12$ in the variance has a simple meaning: each prime contributes $\frac{\cos(t\log p)}{p^{\sigma_0}}$, and the mean of $\cos^2$ is $\frac12$. In the language of Section~\ref{ch:bohrjessen}, $\operatorname{Re}P(\sigma_0+it)$ behaves like $\sum_{p\le X}\frac{\operatorname{Re}X_p}{p^{\sigma_0}}$ with independent uniform $X_p$ on the circle, and the moment computation is a proof that the numbers $p^{-it}$, $p\le X$, have the same joint moments as independent $X_p$ up to order $k$, as long as $X^{2k}=o(T)$.

\section{The mollifier}

The prime sum is Gaussian; we must now connect it to $\zeta$. The plan is to find a Dirichlet polynomial $M$, short enough to be handled by mean values, with $M\approx e^{-P}$ (this section) and $\zeta M\approx1$ (the next). Since $e^{-P}$ is a truncated version of $\prod_p(1-p^{-s})=\sum\mu(n)n^{-s}$, the natural candidate is a truncation of $\sum\mu(n)n^{-s}$. The difficulty is the length. Expanding $e^{-P}$ when $|P|$ is of size $\sqrt{L_2}$ needs more than $\sqrt{L_2}$ terms of the exponential series, and products of that many primes up to $X$ are much larger than $T$. The way out is to split the primes at $Y$: below $Y$, many factors are allowed but the primes are small; between $Y$ and $X$, the primes are large, but $P$ is small there (its variance is only about $2L_3$), so few factors suffice.

Let $K_1=100L_2$, $K_2=100L_3$, and
\[
P_1(s)=\sum_{2\le n\le Y}\frac{\Lambda(n)}{n^s\log n},\qquad P_2(s)=\sum_{Y<n\le X}\frac{\Lambda(n)}{n^s\log n},
\]
so $P=P_1+P_2$. Let $a_1(n)=1$ if $n$ is squarefree, has all its prime factors $\le Y$ and $\omega(n)\le K_1$, and $a_1(n)=0$ otherwise; let $a_2(n)=1$ if $n$ is squarefree, has all its prime factors in $\left]Y,X\right]$ and $\omega(n)\le K_2$, and $0$ otherwise. Put $M_i(s)=\sum_n\mu(n)a_i(n)n^{-s}$ and $M=M_1M_2$. Writing $n=n_1n_2$ with $n_1$ composed of primes $\le Y$ and $n_2$ of primes $>Y$, $M(s)=\sum_n\mu(n)a(n)n^{-s}$ with $a(n)=a_1(n_1)a_2(n_2)$. The mollifier has length
\[
N_M=Y^{K_1}X^{K_2}=T^{100/L_2+100/L_3}=T^{o(1)}.
\]

\begin{lemma}\label{lem:truncexp}
Let $K\ge1$ and $|z|\le K$. Then $\bigl|e^z-\sum_{0\le k\le100K}\frac{z^k}{k!}\bigr|\le e^{-300K}|e^z|$.
\end{lemma}
\begin{proof}
Let $N=\lfloor100K\rfloor+1$. The terms $\frac{K^k}{k!}$ for $k\ge N$ decrease at least by a factor $\frac KN\le\frac1{100}$ at each step, so the tail is at most $\frac{100}{99}\frac{K^N}{N!}\le\frac{100}{99}\bigl(\frac{eK}N\bigr)^N\le\frac{100}{99}\bigl(\frac e{100}\bigr)^{100K}\le e^{-350K}$, using $N!\ge(N/e)^N$. And $|e^z|\ge e^{-K}$.
\end{proof}

Let $\widetilde M_1=\sum_{0\le k\le K_1}\frac{(-P_1)^k}{k!}$ and $\widetilde M_2=\sum_{0\le k\le K_2}\frac{(-P_2)^k}{k!}$, two more Dirichlet polynomials.

\begin{lemma}\label{lem:psmall}
For $T$ large, $\int_T^{2T}|P_1(\sigma_0+it)|^2\dd t\le2TL_2$ and $\int_T^{2T}|P_2(\sigma_0+it)|^2\dd t\le3TL_3$. Consequently $P_T(|P_1|>L_2)\le\frac2{L_2}$, $P_T(|P_2|>L_3)\le\frac3{L_3}$, and outside these two events
\[
\widetilde M_1(\sigma_0+it)=e^{-P_1}\bigl(1+O((\log T)^{-300})\bigr),\qquad\widetilde M_2(\sigma_0+it)=e^{-P_2}\bigl(1+O(L_2^{-300})\bigr).
\]
\end{lemma}
\begin{proof}
By Lemma~\ref{lem:mvt2}, the mean square of $P_1$ is $T\sum_{n\le Y}\frac{\Lambda(n)^2}{n^{2\sigma_0}\log^2n}+O(Y^4)$, and the sum is at most $\sum_{p^k\le Y}\frac1{k^2p^k}\le\sum_{p\le Y}\frac1p+\sum_p\frac1{p(p-1)}\le\log\log Y+O(1)\le L_2+O(1)$. For $P_2$ the sum is at most $\sum_{Y<p\le X}\frac1p+\sum_{p}\sum_{k\ge2,\,p^k>Y}p^{-k}=\log\frac{\log X}{\log Y}+o(1)=2L_3-2\log L_3+o(1)$. Chebyshev's inequality gives the probabilities, and Lemma~\ref{lem:truncexp} with $K=L_2$ and $K=L_3$ the approximations.
\end{proof}

The two truncated exponentials are not quite $M_1$ and $M_2$: the prime powers in $P$ and the truncation create a few extra terms, all of them on integers that are rare in the sense of the weight $n^{-2\sigma_0}$.

\begin{lemma}\label{lem:mollcoef}
For $i=1,2$, $\int_T^{2T}|\widetilde M_i(\sigma_0+it)-M_i(\sigma_0+it)|^2\dd t\le\varepsilon_iT$ with $\varepsilon_1=O((\log T)^{-64})$ and $\varepsilon_2=O(L_2^{-62})$.
\end{lemma}
\begin{proof}
Write $\widetilde M_i=\sum b_i(n)n^{-s}$ and $c_i(n)=b_i(n)-\mu(n)a_i(n)$. Let $\mathcal E$ be the set of $n$ divisible by some $p^j$ with $j\ge2$ and $p^j>Y$.

\emph{Support and size.} Every term of $P_i$ is at most $X$ (at most $Y$ for $i=1$), so $b_i(n)=0$ unless $n\le Y^{K_1}$ (resp.\ $X^{K_2}$) and all primes of $n$ are $\le X$; the same holds for $a_i$. The coefficients of $P_i$ are nonnegative and at most those of $\log\zeta=\sum\frac{\Lambda(n)}{\log n}n^{-s}$, so, as formal Dirichlet series, $|b_i(n)|$ is at most the coefficient of $n$ in $e^{P_i}$, which is at most the coefficient of $n$ in $e^{\log\zeta}=\zeta$, that is $1$. So $|c_i(n)|\le2$.

\emph{Where $c_i(n)=0$.} Let $n\notin\mathcal E$ with $\Omega(n)\le K_i$. Each term of $P_i$ has $n\ge2$, so $(-P_i)^k$ has no term $n^{-s}$ for $k>\Omega(n)$, and $b_i(n)$ is the coefficient of $n^{-s}$ in $e^{-P_i}$. Now $e^{-P_i}=\prod_pe^{-\sum_{j\in J_i(p)}p^{-js}/j}$ with $J_1(p)=\{j:p^j\le Y\}$ and $J_2(p)=\{j:Y<p^j\le X\}$, so $b_i(n)=\prod_{p^j\parallel n}e_{i,p}(j)$, where $e_{i,p}(j)$ is the coefficient of $x^j$ in $\exp(-\sum_{j'\in J_i(p)}x^{j'}/j')$. If $J_i(p)\supset\{1,\dots,j\}$ this is the coefficient of $x^j$ in $\exp(\log(1-x))=1-x$, namely $\mu(p^j)$; if $\min J_i(p)>j$ (or $J_i(p)=\varnothing$), it is $0$. For $i=1$: if a prime $p>Y$ divides $n$, $b_1(n)=0=a_1(n)$; otherwise every $p^j\parallel n$ has $p^j\le Y$ (for $j\ge2$ because $n\notin\mathcal E$), so $b_1(n)=\mu(n)=\mu(n)a_1(n)$. For $i=2$: if some $p\le Y$ divides $n$, then $p^j\le Y$ and $\min J_2(p)>j$, so $b_2(n)=0=a_2(n)$; if some $p>X$ divides $n$, again both vanish; otherwise all primes of $n$ lie in $\left]Y,X\right]$ and $n$ is squarefree (as $n\notin\mathcal E$), $1\in J_2(p)$, and $b_2(n)=\mu(n)=\mu(n)a_2(n)$.

So $c_i(n)\ne0$ forces $n\in\mathcal E$, or $\Omega(n)>K_i$ with all primes of $n$ at most $Y$ (for $i=1$) or in $\left]Y,X\right]$ (for $i=2$, when $n\notin\mathcal E$: a prime $p\le Y$ can only enter $b_2$ through a power $p^j>Y$, $j\ge2$).

\emph{The mean square.} By Lemma~\ref{lem:mvt2}, with $N_i$ the length, $\int|\widetilde M_i-M_i|^2\le4T(S_{\mathcal E}+S_i)+8N_i^4$, where $N_i^4\le T^{400/L_3}=o(T^{1/2})$ and
\[
S_{\mathcal E}=\sum_{\substack{n\in\mathcal E\\p\mid n\Rightarrow p\le X}}\frac1n\le\prod_{p\le X}\Bigl(1-\frac1p\Bigr)^{-1}\sum_{\substack{p\le X,\ j\ge2\\p^j>Y}}\frac1{p^j}\le C\log T\cdot\frac6{\sqrt Y},
\]
because the $p\le\sqrt Y$ contribute at most $\frac2Y$ each and the $p>\sqrt Y$ at most $\frac2{p^2}$. For $S_1$, Rankin's trick with $r=e^{2/3}<2$:
\[
S_1=\sum_{\substack{p\mid n\Rightarrow p\le Y\\\Omega(n)>K_1}}\frac1{n^{2\sigma_0}}\le r^{-K_1}\prod_{p\le Y}\Bigl(1-\frac r{p}\Bigr)^{-1}\le C_rr^{-K_1}(\log Y)^r\le C(\log T)^{-200/3+2},
\]
using $-\log(1-\frac rp)\le\frac rp+\frac{C}{p^2}$ for $p\ge3$ and Mertens. In the same way $S_2\le C_rr^{-K_2}\exp\bigl(r\sum_{Y<p\le X}\frac1p\bigr)\le CL_2^{-200/3}(L_2^2)^r\le CL_2^{-62}$.
\end{proof}

\begin{proposition}\label{prop:mollifier}
$P_T\bigl(|M(\sigma_0+it)e^{P(\sigma_0+it)}-1|>L_2^{-20}\bigr)\to0$.
\end{proposition}
\begin{proof}
Discard the events $|P_1|>L_2$, $|P_2|>L_3$, $|\widetilde M_1-M_1|>(\log T)^{-30}$ and $|\widetilde M_2-M_2|>L_2^{-30}$; by Lemmas~\ref{lem:psmall} and~\ref{lem:mollcoef} and Chebyshev, their probabilities tend to $0$. Outside them, $|e^{P_1}|\le e^{L_2}=\log T$ and
\[
M_1e^{P_1}=\widetilde M_1e^{P_1}+(M_1-\widetilde M_1)e^{P_1}=1+O((\log T)^{-300})+O((\log T)^{-29}),
\]
and similarly, with $|e^{P_2}|\le e^{L_3}=L_2$, $M_2e^{P_2}=1+O(L_2^{-29})$. Multiply: $Me^P=M_1e^{P_1}M_2e^{P_2}=1+O(L_2^{-29})$.
\end{proof}

\section{A twisted second moment}

It remains to show that $\zeta M\approx1$ at $\sigma_0$, in mean square. This is the deepest step, and the only one in which the functional equation of $\zeta$ is used: on a line so close to $\frac12$, $\zeta$ is not well approximated by any Dirichlet polynomial of length below $T$, and the functional equation supplies the missing information. The tool is an exact formula for $|\zeta(s)|^2$ in terms of two absolutely convergent integrals, again due to Radziwi\l\l{} and Soundararajan, in the spirit of the approximate functional equation.

For $s\in\C$ and $c>0$ let
\[
I(s)=\frac1{2\pi i}\int_{(c)}\xi(z+s)\xi(z+\overline s)e^{z^2}\frac{\dd z}z,
\]
the integral along the vertical line $\Re z=c$, upwards. It converges absolutely and does not depend on $c>0$: on a vertical strip $|\xi(z+s)|\le e^{C|z|\log|z|}$ for $|z|$ large, while $|e^{z^2}|=e^{(\Re z)^2-(\Im z)^2}$.

\begin{lemma}\label{lem:exactmod}
For every $s$ with $G(s)\ne0$, $|\zeta(s)|^2=\dfrac{I(s)+I(1-s)}{|G(s)|^2}$.
\end{lemma}
\begin{proof}
Move the line of integration from $\Re z=c$ to $\Re z=-c$; the horizontal segments vanish in the limit by the Gaussian decay. The only pole crossed is $z=0$, with residue $\xi(s)\xi(\overline s)=|\xi(s)|^2$. On $\Re z=-c$ substitute $z=-u$: the line becomes $\Re u=c$ traversed downwards, $\frac{\dd z}z=\frac{\dd u}u$, and $\xi(s-u)=\xi(1-s+u)$, $\xi(\overline s-u)=\xi(1-\overline s+u)$ by the functional equation. Since $\overline{1-s}=1-\overline s$, the integral on $\Re z=-c$ is $-I(1-s)$. So $I(s)=|\xi(s)|^2-I(1-s)$, and $|\xi(s)|^2=|G(s)|^2|\zeta(s)|^2$.
\end{proof}

\begin{lemma}\label{lem:gammaratio}
There are absolute $C$, $T_0$ such that for $T\ge T_0$, $T\le t\le2T$, $\frac12\le\sigma\le1$, $s=\sigma+it$, and $|\Re z|\le2$, $|\Im z|\le\frac T4$,
\[
\frac{G(s+z)G(\overline s+z)}{G(s)G(\overline s)}=\Bigl(\frac t{2\pi}\Bigr)^ze^{E},\qquad|E|\le C\frac{1+|z|^2}T.
\]
Moreover $|G(s)|^{-2}\le Ce^{\pi t/2}$.
\end{lemma}
\begin{proof}
Along the segment from $s$ to $s+z$ the imaginary part stays above $\frac34T$, so Stirling applies to $\frac w2$ for $w$ on the segment, and $\frac{\dd}{\dd w}\log G(w)=\frac1w+\frac1{w-1}-\frac12\log\pi+\frac12\frac{\Gamma'}\Gamma(\frac w2)=\frac12\log\frac w{2\pi}+O(\frac1T)$. For $w=s+u$, $|u|\le|z|\le\frac T2$, $\log\frac w{2\pi}=\log\frac s{2\pi}+O(\frac{|u|}T)$. Integrating, $\log G(s+z)-\log G(s)=\frac z2\log\frac s{2\pi}+O(\frac{|z|+|z|^2}T)$, and the same with $\overline s$. Adding, $\frac z2(\log\frac s{2\pi}+\log\frac{\overline s}{2\pi})=z\log\frac{|s|}{2\pi}=z\log\frac t{2\pi}+O(\frac{|z|}{T^2})$. For the last claim, Stirling gives $\log|\Gamma(\frac s2)|=\frac{\sigma-1}2\log\frac t2-\frac{\pi t}4+O(1)$, so $|G(s)|\ge ct^2t^{-1/4}e^{-\pi t/4}$.
\end{proof}

\begin{theorem}[a twisted second moment]\label{thm:twisted}
There are absolute $C$, $T_0$ such that for $T\ge T_0$, $\frac12<\sigma\le0.6$ and positive integers $h,k\le T$, with $A=\frac{(h,k)^2}{hk}$,
\[
\int_T^{2T}\Bigl(\frac hk\Bigr)^{it}|\zeta(\sigma+it)|^2\dd t=\zeta(2\sigma)A^\sigma T+\zeta(2-2\sigma)A^{1-\sigma}\int_T^{2T}\Bigl(\frac t{2\pi}\Bigr)^{1-2\sigma}\dd t+\mathcal R,
\]
with $|\mathcal R|\le C\bigl(\min(h,k)T^{1-\sigma}(\log T)^3+T^{0.6}\bigr)$.
\end{theorem}
\begin{proof}
\emph{Step 1: replacing the Gamma factors.} Let $\eta=\frac1{\log T}$, $c=1-\sigma+\eta$, and $t\in[T,2T]$. In $I(s)$ on the line $\Re z=c$, $\Re(z+s)=1+\eta$, so $|\zeta(z+s)|,|\zeta(z+\overline s)|\le\zeta(1+\eta)\le2\log T$. On $|\Im z|\le\frac T4$, Lemma~\ref{lem:gammaratio} replaces $\frac{G(z+s)G(z+\overline s)}{|G(s)|^2}$ by $(\frac t{2\pi})^z$ with an error at most $|(\frac t{2\pi})^z|\cdot C\frac{1+|z|^2}T$, and $|(\frac t{2\pi})^z|\le eT^{1-\sigma}$. Since $\int e^{c^2-y^2}\frac{1+|z|^2}{|z|}\dd y\le C\log\frac1c\le C\log\log T$ (with $z=c+iy$), this error contributes $O(T^{-\sigma}(\log T)^3)$ to $I(s)/|G(s)|^2$. On $|\Im z|>\frac T4$, $|G(z+s)|\le C(T+|y|)^2$ (as $|\Gamma(x+iy)|\le\Gamma(x)$) and $|G(s)|^{-2}\le Ce^{\pi T}$, and $e^{-y^2}\le e^{-T^2/16}$ kills everything; the same holds for the integrand with $(\frac t{2\pi})^z$. Therefore
\[
\frac{I(s)}{|G(s)|^2}=\frac1{2\pi i}\int_{(c)}\Bigl(\frac t{2\pi}\Bigr)^z\zeta(z+s)\zeta(z+\overline s)e^{z^2}\frac{\dd z}z+O\bigl(T^{-\sigma}(\log T)^3\bigr).
\]
For $I(1-s)$ take $c'=\sigma+\eta$; since $1-s=\overline s+(1-2\sigma)$ and $1-\overline s=s+(1-2\sigma)$, Lemma~\ref{lem:gammaratio} with $z'=z+1-2\sigma$ (whose real part is $c$) gives in the same way
\[
\frac{I(1-s)}{|G(s)|^2}=\frac1{2\pi i}\int_{(c')}\Bigl(\frac t{2\pi}\Bigr)^{z+1-2\sigma}\zeta(z+1-s)\zeta(z+1-\overline s)e^{z^2}\frac{\dd z}z+O\bigl(T^{-\sigma}(\log T)^3\bigr).
\]

\emph{Step 2: integrating over $t$.} Multiply the first formula by $(h/k)^{it}$, integrate over $[T,2T]$, and expand $\zeta(z+s)\zeta(z+\overline s)=\sum_{m,n}(mn)^{-z-\sigma}(n/m)^{it}$; everything converges absolutely. With $J(u)=\int_T^{2T}(\frac t{2\pi})^u\dd t$, the terms with $hn=km$, that is $m=\frac{Nh}{(h,k)}$, $n=\frac{Nk}{(h,k)}$, $N\ge1$, give
\[
\frac1{2\pi i}\int_{(c)}\zeta(2z+2\sigma)A^{z+\sigma}J(z)e^{z^2}\frac{\dd z}z.
\]
For $hn\ne km$, let $\theta=\log\frac{hn}{km}$. Integrating by parts, $\bigl|\int_T^{2T}e^{i\theta t}(\frac t{2\pi})^z\dd t\bigr|\le(2+|z|)eT^{1-\sigma}\min(T,\frac1{|\theta|})$. So the off-diagonal terms contribute at most $CT^{1-\sigma}\log\log T\cdot S$, with
\[
S=\sum_{hn\ne km}\frac1{(mn)^{1+\eta}}\min\Bigl(T,\frac1{|\log\frac{hn}{km}|}\Bigr).
\]
We show $S\le C\min(h,k)(\log T)^2$. The terms with $\frac{hn}{km}\notin[\frac12,2]$ give at most $\frac{\zeta(1+\eta)^2}{\log2}\le C(\log T)^2$. Fix $m$ and consider the $n$ with $\frac{km}2\le hn\le2km$, $hn\ne km$. Then $\frac1n\le\frac{2h}{km}$, and with $\ell=hn-km\ne0$, $|\log\frac{hn}{km}|\ge\frac{|\ell|}{2km}$. Distinct $n$ give values of $\ell$ spaced by $h$, so for each $j\ge0$ at most two $n$ have $jh\le|\ell|<(j+1)h$; for $j=0$ use $|\ell|\ge1$. Hence the sum over these $n$ is at most
\[
\frac{2h}{km}\cdot2\Bigl(2km+\sum_{1\le j\le km}\frac{2km}{jh}\Bigr)\le8h+8(1+\log km),
\]
and summing over $m$ with the weight $m^{-1-\eta}$ gives at most $C(h\log T+(\log T)^2)$, since $k\le T$ and $\sum_mm^{-1-\eta}\log m\le C(\log T)^2$. Exchanging the roles of $(h,n)$ and $(k,m)$ gives the bound with $k$. The second formula is treated in the same way: the diagonal now gives
\[
\frac1{2\pi i}\int_{(c')}\zeta(2z+2-2\sigma)A^{z+1-\sigma}J(z+1-2\sigma)e^{z^2}\frac{\dd z}z,
\]
and the off-diagonal the same bound, since $|(\frac t{2\pi})^{z+1-2\sigma}|\le eT^{1-\sigma}$ on $\Re z=c'$. With the errors of Step~1, integrated over $[T,2T]$, we are left with the two diagonal integrals, up to $O(\min(h,k)T^{1-\sigma}(\log T)^3)$.

\emph{Step 3: the two main terms.} Put $w=z+\sigma$ in the first diagonal integral and $w=z+1-\sigma$ in the second. Both run over $\Re w=1+\eta$, and their sum is
\[
\frac1{2\pi i}\int_{(1+\eta)}\zeta(2w)A^wJ(w-\sigma)H(w)\dd w,\qquad H(w)=\frac{e^{(w-\sigma)^2}}{w-\sigma}+\frac{e^{(w-1+\sigma)^2}}{w-1+\sigma}.
\]
Move the line to $\Re w=0.1$. On the strip, $|\zeta(2w)|\le C(1+|\Im w|)$ away from $w=\frac12$ (reminder), $|A^w|\le1$, $|J(w-\sigma)|\le T^{2.1}$, and $H$ decays like $e^{-(\Im w)^2}$, so the horizontal segments vanish. The poles crossed are $w=\sigma$, with residue $\zeta(2\sigma)A^\sigma J(0)=\zeta(2\sigma)A^\sigma T$; $w=1-\sigma$, with residue $\zeta(2-2\sigma)A^{1-\sigma}J(1-2\sigma)$; and $w=\frac12$, where $\zeta(2w)$ has a simple pole but $H(\frac12)=\frac{e^{(\sigma-1/2)^2}}{\frac12-\sigma}+\frac{e^{(\sigma-1/2)^2}}{\sigma-\frac12}=0$, so there is no pole. On $\Re w=0.1$, $|J(w-\sigma)|\le T(T/\pi)^{0.1-\sigma}\le T^{0.6}$, $|H(w)|\le Ce^{-(\Im w)^2}$ (as $|w-\sigma|\ge0.4$ and $|w-1+\sigma|\ge0.3$), so the remaining integral is $O(T^{0.6})$.
\end{proof}

All the constants above are uniform in $\sigma\in\left]\frac12,0.6\right]$. Chapter~\ref{ch:moments} lets $\sigma\to\frac12$ and recovers the mean square of $\zeta$ on the critical line.

\begin{lemma}\label{lem:rankinquad}
Let $\mathcal P$ be a finite set of primes, $K\ge0$, and for $p\in\mathcal P$ let $\alpha_p,\beta_p,\gamma_p\in[0,\frac1p]$. For squarefree $h,k$ with prime factors in $\mathcal P$ let $g(h,k)=\prod_{p\mid h,p\nmid k}\alpha_p\prod_{p\mid k,p\nmid h}\beta_p\prod_{p\mid(h,k)}\gamma_p$. Then
\[
\Bigl|\sum_{\omega(h),\omega(k)\le K}\mu(h)\mu(k)g(h,k)-\prod_{p\in\mathcal P}(1-\alpha_p-\beta_p+\gamma_p)\Bigr|\le2e^{-K}\prod_{p\in\mathcal P}\Bigl(1+\frac{2e+1}p\Bigr).
\]
\end{lemma}
\begin{proof}
Without the condition on $\omega$, the sum is multiplicative over the primes of $\mathcal P$, each contributing $1-\alpha_p-\beta_p+\gamma_p$ (according as $p$ divides neither, $h$ only, $k$ only, or both). The pairs with $\omega(h)>K$ contribute at most $\sum e^{\omega(h)-K}g(h,k)=e^{-K}\prod_p(1+e\alpha_p+\beta_p+e\gamma_p)$, and similarly for $k$.
\end{proof}

\begin{proposition}\label{prop:mollmeansq}
$\displaystyle\int_T^{2T}\bigl|1-\zeta(\sigma_0+it)M(\sigma_0+it)\bigr|^2\dd t=o(T)$.
\end{proposition}
\begin{proof}
Write $s=\sigma_0+it$ and expand the square: we need $\int\zeta M=T+o(T)$ and $\int|\zeta M|^2=T+o(T)$. Note $|M(s)|\le N_M$.

\emph{The first moment.} By Volume~14, Lemma~1.5, with $k=3$ and $N=\lceil(3T+3)^{3/2}\rceil$, $\zeta(s)=\sum_{n\le N}n^{-s}+O(T^{-1/4})$ on $[T,2T]$. The integral of $\sum_{n\le N}n^{-s}M(s)$ is $T$ (the term $nm=1$) plus off-diagonal terms with $nm\ge2$, each bounded by $\frac{2}{\log 2}(nm)^{-\sigma_0}$, in total $O(N^{1/2}N_M)=O(T^{3/4+o(1)})$. The error term contributes $O(T^{3/4}N_M)$.

\emph{The second moment.} Expanding,
\[
|\zeta M|^2=|\zeta|^2\sum_{h,k}\frac{\mu(h)\mu(k)a(h)a(k)}{(hk)^{\sigma_0}}\Bigl(\frac kh\Bigr)^{it},
\]
and Theorem~\ref{thm:twisted} applies to each pair. The error terms add up to $O(N_M^3(T^{1/2}(\log T)^3+T^{0.6}))=o(T)$. Since $(hk)^{-\sigma_0}A^{\sigma_0}=\frac{(h,k)^{2\sigma_0}}{(hk)^{2\sigma_0}}$ and $(hk)^{-\sigma_0}A^{1-\sigma_0}=\frac{(h,k)^{2-2\sigma_0}}{hk}$, the main terms are $T\zeta(2\sigma_0)\Sigma+\zeta(2-2\sigma_0)J(1-2\sigma_0)\Sigma'$, with
\[
\Sigma=\sum_{h,k}\mu(h)\mu(k)a(h)a(k)\frac{(h,k)^{2\sigma_0}}{(hk)^{2\sigma_0}},\qquad\Sigma'=\sum_{h,k}\mu(h)\mu(k)a(h)a(k)\frac{(h,k)^{2-2\sigma_0}}{hk}.
\]
Both factor as a sum over the parts with primes $\le Y$ times a sum over the parts with primes in $\left]Y,X\right]$, and Lemma~\ref{lem:rankinquad} applies to each factor, with $K=K_1$ or $K_2$. For $\Sigma$, $\alpha_p=\beta_p=\gamma_p=p^{-2\sigma_0}$; for $\Sigma'$, $\alpha_p=\beta_p=\frac1p$ and $\gamma_p=p^{-2\sigma_0}$. The error terms are $2e^{-K_1}\prod_{p\le Y}(1+\frac{2e+1}p)\le C(\log T)^{-100}(\log T)^{2e+1}\le C(\log T)^{-93}$ and $2e^{-K_2}\prod_{Y<p\le X}(1+\frac{2e+1}p)\le CL_2^{-100}(L_2^2)^{2e+1}\le CL_2^{-87}$.

Let $\delta=2\sigma_0-1=\frac{2W}{\log T}$; then $\delta\log X=2L_3^2\to\infty$ and $\delta\log Y=\frac{2L_3^4}{L_2^2}\to0$. Write $\Pi_1=\prod_{p\le Y}(1-p^{-2\sigma_0})$, $\Pi_2=\prod_{Y<p\le X}(1-p^{-2\sigma_0})$. By the Euler product for $\zeta(2\sigma_0)$, $-\log(1-x)\le2x$ for $x\le\frac12$, and Lemma~\ref{lem:dampedmertens},
\[
1\le\zeta(2\sigma_0)\Pi_1\Pi_2=\prod_{p>X}(1-p^{-2\sigma_0})^{-1}\le\exp\Bigl(2E_1(2L_3^2)+\frac{2C}{\log X}\Bigr)\longrightarrow1,
\]
and likewise $\zeta(2\sigma_0)\Pi_1=\prod_{p>Y}(1-p^{-2\sigma_0})^{-1}\le C(\delta\log Y)^{-2}\le CL_2^4$. With $\zeta(2\sigma_0)\le1+\frac1\delta\le\log T$ and $0\le\Pi_2\le1$,
\[
\zeta(2\sigma_0)\Sigma=\zeta(2\sigma_0)\bigl(\Pi_1+O((\log T)^{-93})\bigr)\bigl(\Pi_2+O(L_2^{-87})\bigr)=1+o(1)+O(L_2^{4-87})+O((\log T)^{-92}).
\]
For the second term, the local factors $1-\frac2p+p^{-2\sigma_0}$ lie in $[0,1-\frac1p]$, so the main products are $\Pi_1'\le\prod_{p\le Y}(1-\frac1p)\le\frac C{\log Y}=\frac{CL_2^2}{\log T}$ and $0\le\Pi_2'\le1$. By the reminder, $|\zeta(2-2\sigma_0)|=|\zeta(1-\delta)|\le\frac1\delta+4\le\log T$, and $J(1-2\sigma_0)\le T(\frac T{2\pi})^{-\delta}\le Te^{-W}$. Hence
\[
\bigl|\zeta(2-2\sigma_0)J(1-2\sigma_0)\Sigma'\bigr|\le Te^{-W}\log T\Bigl(\frac{CL_2^2}{\log T}+C(\log T)^{-93}\Bigr)(1+CL_2^{-87})\le CTL_2^2e^{-L_3^4}=o(T).
\]
So $\int|\zeta M|^2=T+o(T)$, and $\int|1-\zeta M|^2=\int|\zeta M|^2-2\operatorname{Re}\int\zeta M+T=o(T)$.
\end{proof}

The last display is where the two requirements on $W$ meet: $e^{-W}$ must beat $L_2^2$, which comes from $\frac{\log T}{\log Y}$. And $\zeta(2\sigma_0)\Pi_1\Pi_2\to1$ needed $\delta\log X\to\infty$, that is, the mollifier must reach primes beyond $T^{1/W}$, where the damping $p^{-\delta}$ is felt.

\section{Assembling the proof}

\begin{proof}[Proof of Theorem~\ref{thm:selberg}]
Fix $\varepsilon\in\left]0,\frac12\right]$. Outside a set of $t\in[T,2T]$ of measure $o(T)$, three things hold simultaneously:
\begin{enumerate}
\item $|\log|\zeta(\frac12+it)|-\log|\zeta(\sigma_0+it)||\le\varepsilon\sqrt{L_2}$ (Corollary~\ref{cor:offline});
\item $|\zeta(\sigma_0+it)M(\sigma_0+it)-1|\le\varepsilon$ (Proposition~\ref{prop:mollmeansq} and Chebyshev);
\item $|M(\sigma_0+it)e^{P(\sigma_0+it)}-1|\le L_2^{-20}\le\varepsilon$ (Proposition~\ref{prop:mollifier}).
\end{enumerate}
By (2) and (3), $\zeta(\sigma_0+it)\ne0$, $M(\sigma_0+it)\ne0$, and since $|\log|1+x||\le2|x|$ for $|x|\le\frac12$,
\[
\log|\zeta(\sigma_0+it)|=\log|\zeta M|-\log|Me^P|+\operatorname{Re}P=\operatorname{Re}P(\sigma_0+it)+O(\varepsilon).
\]
With (1), $\bigl|\log|\zeta(\frac12+it)|-\operatorname{Re}P(\sigma_0+it)\bigr|\le2\varepsilon\sqrt{L_2}$ outside a set of measure $o(T)$; $\varepsilon$ being arbitrary, the difference divided by $\sqrt{L_2}$ tends to $0$ in probability. By Proposition~\ref{prop:primesgauss}, $\operatorname{Re}P(\sigma_0+it)/\sqrt{v_T/2}\to\mathcal N(0,1)$, and $\frac{v_T}{L_2}\to1$. Slutsky's lemma gives the convergence in law of $\log|\zeta(\frac12+it)|/\sqrt{\frac12L_2}$ to $\mathcal N(0,1)$, which is the theorem since $\Phi$ is continuous.
\end{proof}

\begin{figure}[ht]
\centering\fig{ch07-window}
\caption{The proof at work, on the window $10^6\le t\le10^6+12$. Grey: $\log|\zeta(\frac12+it)|$, plunging towards $-\infty$ at the $22$ zeros of the window (clipped at $-4.5$). Charcoal: $\log|\zeta(\sigma_0+it)|$ with $\sigma_0=\frac12+\frac1{\log T}=0.572$, already smooth. Rust: the prime sum $\operatorname{Re}P(\sigma_0+it)$ over $n\le10^3$, which follows the charcoal curve closely (root mean square difference $0.28$ on the window). On $1500$ random points of $[10^6,2\cdot10^6]$ the correlation between charcoal and rust is $0.95$ with $n\le10^3$ and $0.98$ with $n\le10^5$; between grey and charcoal it is $0.89$.}
\label{fig:selbergwindow}
\end{figure}

\begin{beauty}
Look at Figure~\ref{fig:selbergwindow}. The grey curve has a zero, a point where $\log|\zeta|=-\infty$, every time $t$ moves by about $\frac{2\pi}{\log t}$; the rust one knows nothing of zeros, it is a finite sum of cosines with frequencies $\log p$. Yet in law they are the same. The zeros are where the action is for the primes in short intervals of Volume~14; for the values of $\zeta$ they are, at the scale $\sqrt{\log\log T}$, invisible. And the theorem is the meeting point of two chapters of this book: it is the Erd\H{o}s--Kac theorem with the indicator of $p\mid n$ replaced by the wave $\cos(t\log p)$, and it is the Bohr--Jessen law of Section~\ref{ch:bohrjessen}, in this same chapter, pushed to the edge of the strip, where the limit law dissolves into the Gaussian that it was always approaching.
\end{beauty}

\section{Speed, and what lies beyond}

Figure~\ref{fig:selberghist} shows how slowly the theorem sets in: at $T=10^{10}$ the sample variance is $2.20$ while $\frac12\log\log T=1.57$, and at $T=10^6$ it is $2.00$ against $1.31$. A difference of $0.6$ is not a fluctuation; it is a constant hidden in the $O(1)$ of the variance. Proposition~\ref{prop:primesgauss} computes the variance of the prime sum as $\frac12v_T$, and $v_T=\log\log T+O(\log W)$: the theorem only sees the leading term. Random matrix theory predicts the constant. If $U$ is a Haar unitary matrix of size $N$, Volume~8, Theorem~13.11, shows that $\log|\det(I-U)|$ is Gaussian with variance $\sim\frac12\log N$; the exact variance is $\frac12(\log N+\gamma+1)+o(1)$. With $N=\log\frac t{2\pi}$, the mean density of zeros, this is $2.05$ at $t=1.5\cdot10^6$ and $2.32$ at $t=1.5\cdot10^{10}$, in good agreement with the data. (Keating and Snaith, who made this comparison, also predicted a small arithmetic correction.)

\begin{remark}
These are quoted for culture; see the Sources. (a) Selberg proved the same Gaussian law for $\operatorname{Im}\log\zeta(\frac12+it)=\pi S(t)$, the error term in $N(t)$; the method of this chapter does not give it directly, because $|\zeta M|\approx1$ says nothing about the argument, which could differ from that of $e^P$ by a multiple of $2\pi$. (b) For fixed $\alpha\ne0$, $\log|\zeta(\frac12+it)|$ and $\log|\zeta(\frac12+it+i\alpha)|$ are asymptotically independent Gaussians (Exercise~\ref{ex:twopoint} for the prime sums); Bourgade's multidimensional version, for shifts that shrink with $T$, is not proved in this volume. (c) The rate of convergence: Selberg's original method gives an error $O(\frac{(\log\log\log T)^2}{\sqrt{\log\log T}})$ in the Kolmogorov distance; Roberts (2022) recovered this rate with the method of this chapter. (d) The tails: Radziwi\l\l{} showed that the Gaussian approximation persists for deviations $V\sqrt{\frac12\log\log T}$ with $V$ as large as $(\log\log T)^{1/10-\varepsilon}$, and on the Riemann hypothesis Soundararajan's method gives upper bounds of Gaussian shape for the tail in a much wider range (Chapter~\ref{ch:moments}).
\end{remark}

\section{Exercises}

\begin{exercise}[$\star$]\label{ex:typical}
Deduce from Theorem~\ref{thm:selberg} that for every fixed $C>1$, $P_T\bigl(\frac1C\le|\zeta(\frac12+it)|\le C\bigr)\to0$, and $P_T(|\zeta(\frac12+it)|\le C)\to\frac12$. So $|\zeta(\frac12+it)|$ is typically huge or tiny, each with probability $\frac12$.
\end{exercise}

\begin{exercise}[$\star$]\label{ex:half}
(a) Show that $2^{-k}\binom k{k/2}(k/2)!=2^{-k/2}(k-1)!!$ for even $k$. (b) Let $X_p$ be independent and uniform on the unit circle. Show that $\sum_{p\le X}\frac{\operatorname{Re}X_p}{p^{\sigma_0}}$ has variance $\frac12v_T$, and that its moments of order $k$ are those found in the proof of Proposition~\ref{prop:primesgauss}, up to $O_k(v_T^{k/2-2})$.
\end{exercise}

\begin{exercise}[$\star\star$]\label{ex:complexgauss}
Show that $\bigl(\operatorname{Re}P_0(\sigma_0+it),\operatorname{Im}P_0(\sigma_0+it)\bigr)/\sqrt{v_T/2}$ converges in law to a pair of independent standard Gaussians. Why does this not prove a central limit theorem for $\arg\zeta(\frac12+it)$?
\end{exercise}

\begin{exercise}[$\star\star\star$]\label{ex:twopoint}
Let $\alpha\ne0$ be fixed. Using the prime number theorem in the form $\pi(u)=\int_2^u\frac{\dd v}{\log v}+O(\frac u{\log^2u})$, show that $\sum_{p\le X}p^{-2\sigma_0-i\alpha}=O_\alpha(1)$. Deduce that the vector $\bigl(\operatorname{Re}P_0(\sigma_0+it),\operatorname{Re}P_0(\sigma_0+it+i\alpha)\bigr)/\sqrt{v_T/2}$ converges in law to two independent standard Gaussians.
\end{exercise}

\begin{exercise}[$\star\star\star$]\label{ex:meandiff}
Assume the Riemann hypothesis, and let $\frac12<\sigma\le1$. Show that
\[
\int_T^{2T}\bigl(\log|\zeta(\sigma+it)|-\log|\zeta(\tfrac12+it)|\bigr)\dd t=O\Bigl(\bigl(\sigma-\tfrac12\bigr)\log^2T\Bigr).
\]
In the proof of Proposition~\ref{prop:offline}, the zeros and the Gamma factor each contribute about $\frac12(\sigma-\frac12)T\log T$: show that they cancel.
\end{exercise}

\begin{exercise}[$\star$]\label{ex:fixedW}
Let $\sigma_T=\frac12+\frac{c_T}{\log T}$ with $c_T\ge0$. Show that if $c_T=o(\sqrt{\log\log T})$, then $\log|\zeta(\sigma_T+it)|/\sqrt{\frac12\log\log T}\to\mathcal N(0,1)$ in law. What does Exercise~\ref{ex:varbj} suggest for the variance when $c_T=(\log T)^{1/2}$?
\end{exercise}

\section{Hints}

\begin{hint}{ex:typical}
$\log C/\sqrt{\frac12\log\log T}\to0$; use the continuity of $\Phi$ at $0$.
\end{hint}
\begin{hint}{ex:half}
(a) $(k-1)!!=\frac{k!}{2^{k/2}(k/2)!}$. (b) $\mathbb E[X_p^aX_q^{-b}]$ is $1$ if $p=q$ and $a=b$, and $0$ otherwise.
\end{hint}
\begin{hint}{ex:complexgauss}
Apply the moment computation to $\operatorname{Re}(cP_0)$ for complex $c$, and use characteristic functions in the plane.
\end{hint}
\begin{hint}{ex:twopoint}
Partial summation against $\pi(u)$; the main term is $\int\frac{u^{-\delta-i\alpha}}{u\log u}\dd u$, an oscillating integral in $\log u$. For the vector, the joint moments are sums over matchings, and a pair of primes matched across the two coordinates carries the weight $p^{-2\sigma_0\pm i\alpha}$.
\end{hint}
\begin{hint}{ex:meandiff}
Under the Riemann hypothesis $f_\rho\ge0$ and $\int_\R f_\rho=\pi(\sigma-\frac12)$; count the zeros in $[T,2T]$ with Riemann--von Mangoldt and compare with $\int_T^{2T}\frac12\log\frac t{2\pi}\dd t$.
\end{hint}
\begin{hint}{ex:fixedW}
Proposition~\ref{prop:offline} and Slutsky. For the second question, replace $\sigma-\frac12$ by $\frac{c_T}{\log T}$ in $\frac12\log\frac1{\sigma-\frac12}$.
\end{hint}

\section{Solutions}

\begin{solution}{ex:typical}
\idea Both events say that $\log|\zeta(\frac12+it)|$ is in a window that shrinks to a point after normalisation.
\writeup Let $Y_T=\log|\zeta(\frac12+it)|/\sqrt{\frac12\log\log T}$ and $b_T=\log C/\sqrt{\frac12\log\log T}\to0$. Given $\eta>0$, for $T$ large $b_T\le\eta$, so $P_T(\frac1C\le|\zeta|\le C)=P_T(|Y_T|\le b_T)\le P_T(-\eta\le Y_T\le\eta)\to\Phi(\eta)-\Phi(-\eta)$, which is small with $\eta$. Similarly $P_T(Y_T\le-\eta)\le P_T(|\zeta|\le C)\le P_T(Y_T\le\eta)$, and both bounds tend to $\Phi(\mp\eta)$, which are close to $\frac12$.
\end{solution}

\begin{solution}{ex:half}
\idea Count pairings.
\writeup (a) $2^{-k}\binom k{k/2}(k/2)!=2^{-k}\frac{k!}{(k/2)!}=2^{-k/2}\frac{k!}{2^{k/2}(k/2)!}=2^{-k/2}(k-1)!!$. (b) $\operatorname{Re}X_p=\frac12(X_p+\overline{X_p})$ and $\mathbb E(\operatorname{Re}X_p)^2=\frac12$, so the variance is $\frac12\sum_{p\le X}p^{-2\sigma_0}=\frac12v_T$. For the moments, $\mathbb E\bigl[(\sum_pX_pp^{-\sigma_0})^j(\sum_p\overline{X_p}p^{-\sigma_0})^\ell\bigr]=\sum a_j(n)a_\ell(n)n^{-2\sigma_0}$ exactly, by the orthogonality in the hint and unique factorisation: these are the diagonal terms of the proof, and the off-diagonal terms, which vanish for the model, are the $O(X^{2(j+\ell)}/T)$ there. So the two sides have the same moments up to that error, and the rest of the computation is identical.
\end{solution}

\begin{solution}{ex:complexgauss}
\idea A complex Gaussian is characterised by $\mathbb EZ^j\overline Z^\ell=\delta_{j\ell}\,j!$.
\writeup For $c\in\C$,
\[
(\operatorname{Re}cP_0)^k=2^{-k}\sum_j\binom kjc^j\overline c^{k-j}P_0^j\overline{P_0}^{k-j},
\]
and by the proof of Proposition~\ref{prop:primesgauss} its mean is $2^{-k}\binom k{k/2}|c|^k(k/2)!v_T^{k/2}+O_k(v_T^{k/2-2})+o(1)$ for even $k$, and $o(1)$ for odd $k$. So $\operatorname{Re}(cP_0)/\sqrt{v_T/2}\to\mathcal N(0,|c|^2)$. With $c=a-ib$, $\operatorname{Re}(cP_0)=a\operatorname{Re}P_0+b\operatorname{Im}P_0$, and $a^2+b^2=|c|^2$ is the variance of $aU+bV$ for independent standard Gaussians $U,V$. Hence the characteristic function of the normalised pair at $(a,b)$ converges to $e^{-(a^2+b^2)/2}$, and Lévy's continuity theorem in $\R^2$ concludes. For the argument of $\zeta$, the chain of approximations only controls $|\zeta M|$: the step $\zeta M\approx1$ in mean square does give $\arg(\zeta M)\approx0$ modulo $2\pi$, but $\arg\zeta(\frac12+it)$ is defined by continuous variation along a path, and the passage from $\sigma_0$ to $\frac12$ (Proposition~\ref{prop:offline}) has no analogue for the argument: near a zero, $\arg\zeta$ jumps by $\pi$ on a very short interval, and the average control of Lemma~\ref{lem:logint} is lost.
\end{solution}

\begin{solution}{ex:twopoint}
\idea The covariance between the two coordinates is a sum over primes that oscillates, so it stays bounded while the variances grow like $\log\log T$.
\writeup Let $\delta=2\sigma_0-1$ and $R(u)=\pi(u)-\operatorname{li}(u)$, with $\operatorname{li}(u)=\int_2^u\frac{\dd v}{\log v}$. Then
\[
\sum_{p\le X}p^{-1-\delta-i\alpha}=\int_2^X\frac{u^{-1-\delta-i\alpha}}{\log u}\dd u+\int_{2^-}^Xu^{-1-\delta-i\alpha}\dd R(u).
\]
In the first integral put $u=e^v$: it is $\int_{\log2}^{\log X}e^{-(\delta+i\alpha)v}\frac{\dd v}v$, and integrating by parts once, with $e^{-(\delta+i\alpha)v}=\frac{\dd}{\dd v}\frac{e^{-(\delta+i\alpha)v}}{-(\delta+i\alpha)}$, it is bounded by $\frac{C}{|\alpha|}$. The second, by parts, is $O(\frac1{\log^2X})+|1+\delta+i\alpha|\int_2^\infty\frac{u^{-2}u}{\log^2u}\dd u=O_\alpha(1)$. So $w_T=\sum_{p\le X}p^{-2\sigma_0-i\alpha}=O_\alpha(1)$.

Now $P_0(\sigma_0+it+i\alpha)=\sum_{p\le X}p^{-i\alpha}p^{-\sigma_0-it}$ is again a Dirichlet polynomial in $t$, with coefficients of modulus $p^{-\sigma_0}$. The proof of Proposition~\ref{prop:primesgauss} applies to any product of $j_1$ factors $P_0(s)$, $\ell_1$ factors $\overline{P_0(s)}$, $j_2$ factors $P_0(s+i\alpha)$ and $\ell_2$ factors $\overline{P_0(s+i\alpha)}$: the mean is $o(1)$ plus the diagonal sum, in which the primes of the unconjugated factors must be a permutation of those of the conjugated ones. Up to $O(v_T^{(j_1+\ell_1+j_2+\ell_2)/2-2})$ (repeated primes), the diagonal is a sum over the bijections between unconjugated and conjugated factors, each pair contributing $v_T$ if both factors belong to the same coordinate and $w_T$ or $\overline{w_T}$ otherwise. These are the moments of a complex Gaussian vector with covariance $\begin{pmatrix}v_T&w_T\\\overline{w_T}&v_T\end{pmatrix}$, and after division by $v_T$ the off-diagonal entries tend to $0$. Taking real parts and arguing as in Exercise~\ref{ex:complexgauss} (with $a\operatorname{Re}P_0(s)+b\operatorname{Re}P_0(s+i\alpha)$), the pair converges to independent standard Gaussians.
\end{solution}

\begin{solution}{ex:meandiff}
\idea The mean of $\log|\zeta(\sigma+it)|$ should not depend on $\sigma$ near $\frac12$, and the proof of Proposition~\ref{prop:offline} shows the two large contributions explicitly.
\writeup By the proof of Proposition~\ref{prop:offline}, the integral equals $\int_T^{2T}\sum_\rho f_\rho-\int_T^{2T}\bigl(\log|G(\sigma+it)|-\log|G(\frac12+it)|\bigr)\dd t$. The Gamma part is $(\sigma-\frac12)\int_T^{2T}\frac12\log\frac t{2\pi}\dd t+O(\sigma-\frac12)$, since $\operatorname{Re}\frac{G'}G(u+it)=\frac12\log\frac t{2\pi}+O(\frac1t)$. Under the Riemann hypothesis $\beta=\frac12$, so $f_\rho\ge0$ and $\int_\R f_\rho=\pi(\sigma-\frac12)$ by Lemma~\ref{lem:logint}. A zero with $\gamma\in[T,2T]$ at distance $d$ from the endpoints loses $\int_{\R\setminus[T,2T]}f_\rho\le(\sigma-\frac12)\min(\pi,\frac1d)$, and a zero outside, at distance $d$, contributes at most the same amount. There are $O(\log T)$ zeros at distance in $[j,j+1[$ from each endpoint, for $j\le T$, so these corrections total $O((\sigma-\frac12)\log T\sum_{j\le T}\frac1{j+1})+O((\sigma-\frac12)\sqrt T\log T/\sqrt T)=O((\sigma-\frac12)\log^2T)$, the zeros at distance $>T$ being handled by $|f_\rho|\le\frac{\sigma-\frac12}{d^2}$ as before. So $\int_T^{2T}\sum_\rho f_\rho=\pi(\sigma-\frac12)(N(2T)-N(T))+O((\sigma-\frac12)\log^2T)$. By Riemann--von Mangoldt, $N(2T)-N(T)=\frac1{2\pi}\bigl[u\log\frac u{2\pi e}\bigr]_T^{2T}+O(\log T)=\frac1{2\pi}\int_T^{2T}\log\frac t{2\pi}\dd t+O(\log T)$, because $\frac{\dd}{\dd u}u\log\frac u{2\pi e}=\log\frac u{2\pi}$. The main terms $\frac12(\sigma-\frac12)\int_T^{2T}\log\frac t{2\pi}\dd t$ cancel exactly.
\end{solution}

\begin{solution}{ex:fixedW}
\idea The cost of moving off the line is $(\sigma-\frac12)\log T$ on average, and Selberg's normalisation absorbs anything smaller than $\sqrt{\log\log T}$.
\writeup By Proposition~\ref{prop:offline} and Markov's inequality, $P_T\bigl(|\log|\zeta(\frac12+it)|-\log|\zeta(\sigma_T+it)||>\varepsilon\sqrt{\log\log T}\bigr)\le\frac{Cc_T}{\varepsilon\sqrt{\log\log T}}\to0$ (for $c_T=0$ there is nothing to prove). So the difference of the two normalised variables tends to $0$ in probability, and Theorem~\ref{thm:selberg} and Slutsky's lemma give the claim. For $c_T=(\log T)^{1/2}$, $\sigma_T-\frac12=(\log T)^{-1/2}$, and the Bohr--Jessen variance $\frac12\log\frac1{\sigma_T-\frac12}$ becomes $\frac14\log\log T$: half of Selberg's. Off the line by $(\log T)^{-1/2}$, only the primes up to about $e^{\sqrt{\log T}}$ contribute, and $\log\log$ of that is $\frac12\log\log T$. (The proof of this chapter can be adapted to make this heuristic rigorous; we do not pursue it.)
\end{solution}

\section*{Sources and further reading}

\begin{itemize}
\item \emph{Selberg's war.} The episodes of 1940--1945 (the fighting, the prison camp at Trandum, the arrest after the doctoral examination, the solitary research in Gj\o{}vik) and the Fields Medal of 1950 are from J.~J.~O'Connor and E.~F.~Robertson, ``Atle Selberg'', \emph{MacTutor History of Mathematics Archive}, University of St Andrews.
\item \emph{Selberg's papers.} A.~Selberg, ``On the zeros of Riemann's zeta-function'', \emph{Skr.\ Norske Vid.-Akad.\ Oslo I} 1942, no.~10. A.~Selberg, ``Contributions to the theory of the Riemann zeta-function'', \emph{Arch.\ Math.\ Naturvid.}\ 48 (1946), no.~5, 89--155. A.~Selberg, ``Old and new conjectures and results about a class of Dirichlet series'', in \emph{Proceedings of the Amalfi Conference on Analytic Number Theory} (Maiori, 1989), Univ.\ Salerno, 1992, 367--385.
\item \emph{The proof of this chapter.} M.~Radziwi\l\l{} and K.~Soundararajan, ``Selberg's central limit theorem for $\log|\zeta(\frac12+it)|$'', \emph{L'Enseignement Math.}\ (2) 63 (2017), 1--19 (arXiv:1509.06827). We follow its four propositions, its parameters and its proof of the twisted second moment, with the details written out.
\item \emph{Other proofs.} E.~Bombieri and D.~A.~Hejhal, ``On the distribution of zeros of linear combinations of Euler products'', \emph{Duke Math.\ J.}\ 80 (1995), 821--862. A.~Laurinčikas, ``A limit theorem for the Riemann zeta-function on the critical line. II'', \emph{Litovsk.\ Mat.\ Sb.}\ 27 (1987), 489--500.
\item \emph{Random matrices.} J.~P.~Keating and N.~C.~Snaith, ``Random matrix theory and $\zeta(1/2+it)$'', \emph{Comm.\ Math.\ Phys.}\ 214 (2000), 57--89.
\item \emph{Rate and tails.} A.~Roberts, ``The multivariate rate of convergence for Selberg's central limit theorem'', arXiv:2212.01411 (2022). M.~Radziwi\l\l, ``Large deviations in Selberg's central limit theorem'', arXiv:1108.5092 (2011).
\end{itemize}


\section{Exercises}

\begin{exercise}[$\star$]\label{ex:meansquarezeta}
Show that $\mathbb E|\zeta(\sigma,X)|^2=\zeta(2\sigma)$. Prove, by the method of Proposition~\ref{prop:mollified} with $y=1$ (so $A_1=1$), that $\mathbb E_T|\zeta(\sigma+it)|^2\to\zeta(2\sigma)$: the mean square of $\zeta$ on a vertical line is the mean square of the model. (This is a theorem of Hardy and Littlewood, 1918.)
\end{exercise}

\begin{exercise}[$\star$]\label{ex:sigmabig}
Let $\sigma>1$. Show that $\frac{\zeta(2\sigma)}{\zeta(\sigma)}\le|\zeta(\sigma+it)|\le\zeta(\sigma)$ for every $t$, so the values of $\zeta$ on the line $\Re s=\sigma$ are not dense. Which step of Lemma~\ref{lem:support} fails?
\end{exercise}

\begin{exercise}[$\star\star$]\label{ex:kwfail}
Show that $(2^{-it},4^{-it})$, $t$ uniform on $[T,2T]$, converges in law to $(X,X^2)$ with $X$ uniform on the circle. What happens for $(2^{-it},3^{-it},6^{-it})$? Formulate the general statement for $(n_1^{-it},\dots,n_r^{-it})$.
\end{exercise}

\begin{exercise}[$\star\star$]\label{ex:varbj}
Show that $\log|\zeta(\sigma,X)|$ has mean $0$ and variance $v(\sigma)=\frac12\sum_p\sum_{k\ge1}\frac1{k^2p^{2k\sigma}}$, and that $v(\sigma)=\frac12\log\frac1{\sigma-\frac12}+O(1)$ as $\sigma\to\frac12^+$. Compare with the variance $\frac12\log\log T$ in Selberg's theorem of Section~\ref{ch:selbergclt}, at $\sigma-\frac12=\frac1{\log T}$.
\end{exercise}

\begin{exercise}[$\star\star$]\label{ex:meanzeta}
Show directly, with Volume~14, Lemma~1.5, that $\mathbb E_T\zeta(\sigma+it)\to1=\mathbb E\zeta(\sigma,X)$. Why does this not follow from Theorem~\ref{thm:bohrjessen} alone?
\end{exercise}

\begin{exercise}[$\star\star\star$]\label{ex:jointbj}
Let $\chi$ be the non-principal character modulo $4$. Show that $(\zeta(\sigma+it),L(\sigma+it,\chi))$ converges in law to
\[
\Bigl(\prod_p(1-X_pp^{-\sigma})^{-1},\ \prod_p(1-\chi(p)X_pp^{-\sigma})^{-1}\Bigr),
\]
and deduce that the two values are \emph{not} asymptotically independent. (Assume the analogue of Proposition~\ref{prop:mollified} for $L(s,\chi)$, proved in the same way.)
\end{exercise}

\begin{exercise}[$\star$]\label{ex:typical}
Deduce from Theorem~\ref{thm:selberg} that for every fixed $C>1$, $P_T\bigl(\frac1C\le|\zeta(\frac12+it)|\le C\bigr)\to0$, and $P_T(|\zeta(\frac12+it)|\le C)\to\frac12$. So $|\zeta(\frac12+it)|$ is typically huge or tiny, each with probability $\frac12$.
\end{exercise}

\begin{exercise}[$\star$]\label{ex:half}
(a) Show that $2^{-k}\binom k{k/2}(k/2)!=2^{-k/2}(k-1)!!$ for even $k$. (b) Let $X_p$ be independent and uniform on the unit circle. Show that $\sum_{p\le X}\frac{\operatorname{Re}X_p}{p^{\sigma_0}}$ has variance $\frac12v_T$, and that its moments of order $k$ are those found in the proof of Proposition~\ref{prop:primesgauss}, up to $O_k(v_T^{k/2-2})$.
\end{exercise}

\begin{exercise}[$\star\star$]\label{ex:complexgauss}
Show that $\bigl(\operatorname{Re}P_0(\sigma_0+it),\operatorname{Im}P_0(\sigma_0+it)\bigr)/\sqrt{v_T/2}$ converges in law to a pair of independent standard Gaussians. Why does this not prove a central limit theorem for $\arg\zeta(\frac12+it)$?
\end{exercise}

\begin{exercise}[$\star\star\star$]\label{ex:twopoint}
Let $\alpha\ne0$ be fixed. Using the prime number theorem in the form $\pi(u)=\int_2^u\frac{\dd v}{\log v}+O(\frac u{\log^2u})$, show that $\sum_{p\le X}p^{-2\sigma_0-i\alpha}=O_\alpha(1)$. Deduce that the vector $\bigl(\operatorname{Re}P_0(\sigma_0+it),\operatorname{Re}P_0(\sigma_0+it+i\alpha)\bigr)/\sqrt{v_T/2}$ converges in law to two independent standard Gaussians.
\end{exercise}

\begin{exercise}[$\star\star\star$]\label{ex:meandiff}
Assume the Riemann hypothesis, and let $\frac12<\sigma\le1$. Show that
\[
\int_T^{2T}\bigl(\log|\zeta(\sigma+it)|-\log|\zeta(\tfrac12+it)|\bigr)\dd t=O\Bigl(\bigl(\sigma-\tfrac12\bigr)\log^2T\Bigr).
\]
In the proof of Proposition~\ref{prop:offline}, the zeros and the Gamma factor each contribute about $\frac12(\sigma-\frac12)T\log T$: show that they cancel.
\end{exercise}

\begin{exercise}[$\star$]\label{ex:fixedW}
Let $\sigma_T=\frac12+\frac{c_T}{\log T}$ with $c_T\ge0$. Show that if $c_T=o(\sqrt{\log\log T})$, then $\log|\zeta(\sigma_T+it)|/\sqrt{\frac12\log\log T}\to\mathcal N(0,1)$ in law. What does Exercise~\ref{ex:varbj} suggest for the variance when $c_T=(\log T)^{1/2}$?
\end{exercise}

\section{Hints}

\begin{hint}{ex:meansquarezeta}
For the model, Proposition~\ref{prop:randomeuler} and independence. For $\zeta$, the coefficients are $a_n=n^{-\sigma}$ for $n\le N$.
\end{hint}
\begin{hint}{ex:sigmabig}
$|1-p^{-s}|$ lies between $1-p^{-\sigma}$ and $1+p^{-\sigma}$, and $\prod_p(1+p^{-\sigma})=\frac{\zeta(\sigma)}{\zeta(2\sigma)}$.
\end{hint}
\begin{hint}{ex:kwfail}
Compute the limits of $\mathbb E_T2^{-itm_1}4^{-itm_2}$.
\end{hint}
\begin{hint}{ex:varbj}
Orthogonality of $X_p^k$; then Mertens in the form $\sum_pp^{-1-2\delta}=\log\frac1{2\delta}+O(1)$.
\end{hint}
\begin{hint}{ex:meanzeta}
Integrate $\sum_{n\le N}n^{-s}$ term by term. Convergence in law does not control expectations of unbounded functions.
\end{hint}
\begin{hint}{ex:jointbj}
Kronecker--Weyl for pairs of continuous functions of $(p^{-it})_{p\le y}$; for the dependence, compare $\mathbb E\,\zeta\overline{L}$ with $\mathbb E\zeta\,\mathbb E\overline L$ in the model.
\end{hint}

\begin{hint}{ex:typical}
$\log C/\sqrt{\frac12\log\log T}\to0$; use the continuity of $\Phi$ at $0$.
\end{hint}
\begin{hint}{ex:half}
(a) $(k-1)!!=\frac{k!}{2^{k/2}(k/2)!}$. (b) $\mathbb E[X_p^aX_q^{-b}]$ is $1$ if $p=q$ and $a=b$, and $0$ otherwise.
\end{hint}
\begin{hint}{ex:complexgauss}
Apply the moment computation to $\operatorname{Re}(cP_0)$ for complex $c$, and use characteristic functions in the plane.
\end{hint}
\begin{hint}{ex:twopoint}
Partial summation against $\pi(u)$; the main term is $\int\frac{u^{-\delta-i\alpha}}{u\log u}\dd u$, an oscillating integral in $\log u$. For the vector, the joint moments are sums over matchings, and a pair of primes matched across the two coordinates carries the weight $p^{-2\sigma_0\pm i\alpha}$.
\end{hint}
\begin{hint}{ex:meandiff}
Under the Riemann hypothesis $f_\rho\ge0$ and $\int_\R f_\rho=\pi(\sigma-\frac12)$; count the zeros in $[T,2T]$ with Riemann--von Mangoldt and compare with $\int_T^{2T}\frac12\log\frac t{2\pi}\dd t$.
\end{hint}
\begin{hint}{ex:fixedW}
Proposition~\ref{prop:offline} and Slutsky. For the second question, replace $\sigma-\frac12$ by $\frac{c_T}{\log T}$ in $\frac12\log\frac1{\sigma-\frac12}$.
\end{hint}

\section{Solutions}

\begin{solution}{ex:meansquarezeta}
\idea The model's second moment factors over primes; for $\zeta$, the short sum and the crude mean value theorem show that only the diagonal survives.
\writeup By independence and Proposition~\ref{prop:randomeuler}, $\mathbb E|\zeta_y(X)|^2=\prod_{p\le y}(1-p^{-2\sigma})^{-1}\to\zeta(2\sigma)$, and $|\zeta_y(X)|^2\to|\zeta(\sigma,X)|^2$ almost surely. The convergence of expectations needs uniform integrability: $\zeta_y(X)$ is a martingale in $y$ (multiply by an independent factor of mean $1$, since $\mathbb E(1-X_pp^{-\sigma})^{-1}=1$), bounded in $L^2$ by $\zeta(2\sigma)$, so it converges in $L^2$ (Volume~4, Chapter~10) and $\mathbb E|\zeta(\sigma,X)|^2=\zeta(2\sigma)$. For $\zeta$: as in Proposition~\ref{prop:mollified} with $A_y$ replaced by $1$, $\zeta=Z_N+R$ with $R\to0$ uniformly, and Lemma~\ref{lem:crudemvt} with $a_n=n^{-\sigma}$ gives $\mathbb E_T|Z_N|^2=\sum_{n\le N}n^{-2\sigma}+O(\frac{N^{2-2\sigma}\log N}T)\to\zeta(2\sigma)$. Then $\mathbb E_T|\zeta|^2\to\zeta(2\sigma)$.
\end{solution}

\begin{solution}{ex:sigmabig}
\idea For $\sigma>1$ every Euler factor is pinched between two positive constants, and the product of the bounds converges.
\writeup $|\zeta(s)|=\prod_p|1-p^{-s}|^{-1}$ with $1-p^{-\sigma}\le|1-p^{-s}|\le1+p^{-\sigma}$. So $|\zeta(s)|\le\prod_p(1-p^{-\sigma})^{-1}=\zeta(\sigma)$ and $|\zeta(s)|\ge\prod_p(1+p^{-\sigma})^{-1}=\prod_p\frac{1-p^{-\sigma}}{1-p^{-2\sigma}}=\frac{\zeta(2\sigma)}{\zeta(\sigma)}$. The values lie in a closed annulus, so they are not dense. In Lemma~\ref{lem:support}, the greedy choice needs $\sum_pp^{-\sigma}=\infty$ to travel arbitrarily far; for $\sigma>1$ the total length available is finite.
\end{solution}

\begin{solution}{ex:kwfail}
\idea Characters of the torus see only the products $\prod n_j^{m_j}$, and the limit is uniform on the subgroup cut out by the multiplicative relations.
\writeup $\mathbb E_T2^{-itm_1}4^{-itm_2}=\mathbb E_T2^{-it(m_1+2m_2)}\to\mathbf 1\{m_1+2m_2=0\}$, which is $\mathbb EX^{m_1}(X^2)^{m_2}$. By Stone--Weierstrass, as in Lemma~\ref{lem:kw}, the limit law is that of $(X,X^2)$: concentrated on the curve $z_2=z_1^2$, not uniform on the torus. Similarly $(2^{-it},3^{-it},6^{-it})\Rightarrow(X_2,X_3,X_2X_3)$. In general the limit of $(n_j^{-it})$ is the law of $(f(n_j))_j$ for the Steinhaus random multiplicative function $f$ of Chapter~\ref{ch:randommult}: this is the precise sense in which $n^{-it}$ is a random multiplicative function.
\end{solution}

\begin{solution}{ex:varbj}
\idea Different powers of a uniform point of the circle are orthogonal, so the variance is a sum of squares of Taylor coefficients.
\writeup $\Re\ell_p(X_p)=\Re\sum_ka_kX_p^k$ with $a_k=\frac1{kp^{k\sigma}}$ real. Its mean is $0$, and $\mathbb E(\Re\sum a_kX^k)^2=\frac12\mathbb E|\sum a_kX^k|^2+\frac12\Re\mathbb E(\sum a_kX^k)^2=\frac12\sum a_k^2$, because $\mathbb EX^{j+k}=0$ for $j,k\ge1$. Summing over independent primes (the series converges in $L^2$) gives $v(\sigma)$. Its mean $0$ follows similarly. As $\sigma\to\frac12^+$, the terms $k\ge2$ contribute $O(\sum_pp^{-2})=O(1)$, and $\frac12\sum_pp^{-2\sigma}=\frac12\sum_pp^{-1-2\delta}$ with $\delta=\sigma-\frac12$. By partial summation from Mertens' theorem (Volume~6, Theorem~1.5), $\sum_pp^{-1-2\delta}=\log\frac1{2\delta}+O(1)$ for $0<\delta\le\frac12$. So $v(\sigma)=\frac12\log\frac1{\sigma-\frac12}+O(1)$. At $\sigma-\frac12=\frac1{\log T}$ this is $\frac12\log\log T+O(1)$: exactly Selberg's variance. The random Euler product, pushed to distance $\frac1{\log T}$ from the critical line, already knows the answer of Section~\ref{ch:selbergclt}.
\end{solution}

\begin{solution}{ex:meanzeta}
\idea Every term $n^{-s}$ with $n\ge2$ averages to almost nothing; the term $n=1$ gives $1$.
\writeup With $N$ as in Proposition~\ref{prop:mollified}, $\zeta=Z_N+R$ with $R\to0$ uniformly on $[T,2T]$. And $\mathbb E_TZ_N=1+\sum_{2\le n\le N}n^{-\sigma}\mathbb E_Tn^{-it}$ with $|\mathbb E_Tn^{-it}|\le\frac2{T\log n}$, so the sum is at most $\frac2{T\log2}\sum_{n\le N}n^{-\sigma}\ll\frac{N^{1-\sigma}}T\to0$. For the model, $\mathbb E\zeta_y(X)=\prod_{p\le y}\mathbb E\sum_kX_p^kp^{-k\sigma}=1$, and the $L^2$ convergence of Exercise~\ref{ex:meansquarezeta} gives $\mathbb E\zeta(\sigma,X)=1$. Theorem~\ref{thm:bohrjessen} concerns bounded test functions; $z\mapsto z$ is unbounded, and convergence in law alone does not prevent a small set of $t$ with huge values of $\zeta$ from shifting the mean. Here the bounded second moments (Exercise~\ref{ex:meansquarezeta}) give uniform integrability, which is the missing ingredient.
\end{solution}

\begin{solution}{ex:jointbj}
\idea Both functions are built from the same points $p^{-it}$, so their limits are built from the same $X_p$.
\writeup For fixed $y$, $(E_y(t),E_y^\chi(t))$ is a continuous function of $(p^{-it})_{p\le y}$, so by Lemma~\ref{lem:kw} it converges in law to the pair formed by $\zeta_y(X)$ and $\prod_{p\le y}(1-\chi(p)X_pp^{-\sigma})^{-1}$. The comparison step of Theorem~\ref{thm:bohrjessen}, for $g$ Lipschitz on $\C^2$, uses Proposition~\ref{prop:mollified} for $\zeta$ and its analogue for $L$, and the limit $y\to\infty$ is as before. In the model, since $\chi$ is real,
\[
\mathbb E\zeta(\sigma,X)\overline{L(\sigma,X)}=\prod_p\sum_{k\ge0}\frac{\chi(p)^k}{p^{2k\sigma}}=\prod_p\Bigl(1-\frac{\chi(p)}{p^{2\sigma}}\Bigr)^{-1}=L(2\sigma,\chi),
\]
by orthogonality of the powers of $X_p$, while $\mathbb E\zeta\,\mathbb E\overline L=1$. Since $L(2\sigma,\chi)\ne1$ (for instance it is at most $1-3^{-2\sigma}+5^{-2\sigma}<1$ by the alternating series), the two limits are correlated, hence not independent. (Both have finite second moments, by the argument of Exercise~\ref{ex:meansquarezeta}, so the covariance is defined.)
\end{solution}

\section*{Sources and further reading}

\begin{itemize}
\item \emph{Bohr the footballer.} The 1908 Olympic football tournament and Bohr's silver medal are recorded in the MacTutor biography: J.~J.~O'Connor and E.~F.~Robertson, ``Harald August Bohr'', \emph{MacTutor History of Mathematics Archive}, University of St Andrews.
\item \emph{Density of the values.} H.~Bohr and R.~Courant, ``Neue Anwendungen der Theorie der Diophantischen Approximationen auf die Riemannsche Zetafunktion'', \emph{J.\ Reine Angew.\ Math.}\ 144 (1914), 249--274.
\item \emph{The limit law.} H.~Bohr and B.~Jessen, ``Über die Werteverteilung der Riemannschen Zetafunktion'', \emph{Acta Math.}\ 54 (1930), 1--35, and 58 (1932), 1--55. B.~Jessen and A.~Wintner, ``Distribution functions and the Riemann zeta function'', \emph{Trans.\ Amer.\ Math.\ Soc.}\ 38 (1935), 48--88.
\item \emph{Universality.} S.~M.~Voronin, ``Theorem on the `universality' of the Riemann zeta-function'', \emph{Izv.\ Akad.\ Nauk SSSR Ser.\ Mat.}\ 39 (1975), 475--486. B.~Bagchi, \emph{The statistical behaviour and universality properties of the Riemann zeta-function and other allied Dirichlet series}, PhD thesis, Indian Statistical Institute, Calcutta, 1981.
\item \emph{Tails.} H.~L.~Montgomery, ``Extreme values of the Riemann zeta function'', \emph{Comment.\ Math.\ Helv.}\ 52 (1977), 511--518. Y.~Lamzouri, ``On the distribution of extreme values of zeta and $L$-functions in the strip $\frac12<\sigma<1$'', \emph{Int.\ Math.\ Res.\ Not.}\ 2011, 5449--5503.
\item \emph{Books.} E.~Kowalski, \emph{Arithmetic Randonnée: An Introduction to Probabilistic Number Theory}, Cambridge University Press, 2021, Chapter~3, whose proof of the Bohr--Jessen theorem we follow in spirit. A.~Laurinčikas, \emph{Limit Theorems for the Riemann Zeta-Function}, Kluwer, 1996. E.~C.~Titchmarsh, \emph{The Theory of the Riemann Zeta-Function}, 2nd ed., Oxford, 1986, Chapter~11.
\end{itemize}

\begin{solution}{ex:typical}
\idea Both events say that $\log|\zeta(\frac12+it)|$ is in a window that shrinks to a point after normalisation.
\writeup Let $Y_T=\log|\zeta(\frac12+it)|/\sqrt{\frac12\log\log T}$ and $b_T=\log C/\sqrt{\frac12\log\log T}\to0$. Given $\eta>0$, for $T$ large $b_T\le\eta$, so $P_T(\frac1C\le|\zeta|\le C)=P_T(|Y_T|\le b_T)\le P_T(-\eta\le Y_T\le\eta)\to\Phi(\eta)-\Phi(-\eta)$, which is small with $\eta$. Similarly $P_T(Y_T\le-\eta)\le P_T(|\zeta|\le C)\le P_T(Y_T\le\eta)$, and both bounds tend to $\Phi(\mp\eta)$, which are close to $\frac12$.
\end{solution}

\begin{solution}{ex:half}
\idea Count pairings.
\writeup (a) $2^{-k}\binom k{k/2}(k/2)!=2^{-k}\frac{k!}{(k/2)!}=2^{-k/2}\frac{k!}{2^{k/2}(k/2)!}=2^{-k/2}(k-1)!!$. (b) $\operatorname{Re}X_p=\frac12(X_p+\overline{X_p})$ and $\mathbb E(\operatorname{Re}X_p)^2=\frac12$, so the variance is $\frac12\sum_{p\le X}p^{-2\sigma_0}=\frac12v_T$. For the moments, $\mathbb E\bigl[(\sum_pX_pp^{-\sigma_0})^j(\sum_p\overline{X_p}p^{-\sigma_0})^\ell\bigr]=\sum a_j(n)a_\ell(n)n^{-2\sigma_0}$ exactly, by the orthogonality in the hint and unique factorisation: these are the diagonal terms of the proof, and the off-diagonal terms, which vanish for the model, are the $O(X^{2(j+\ell)}/T)$ there. So the two sides have the same moments up to that error, and the rest of the computation is identical.
\end{solution}

\begin{solution}{ex:complexgauss}
\idea A complex Gaussian is characterised by $\mathbb EZ^j\overline Z^\ell=\delta_{j\ell}\,j!$.
\writeup For $c\in\C$,
\[
(\operatorname{Re}cP_0)^k=2^{-k}\sum_j\binom kjc^j\overline c^{k-j}P_0^j\overline{P_0}^{k-j},
\]
and by the proof of Proposition~\ref{prop:primesgauss} its mean is $2^{-k}\binom k{k/2}|c|^k(k/2)!v_T^{k/2}+O_k(v_T^{k/2-2})+o(1)$ for even $k$, and $o(1)$ for odd $k$. So $\operatorname{Re}(cP_0)/\sqrt{v_T/2}\to\mathcal N(0,|c|^2)$. With $c=a-ib$, $\operatorname{Re}(cP_0)=a\operatorname{Re}P_0+b\operatorname{Im}P_0$, and $a^2+b^2=|c|^2$ is the variance of $aU+bV$ for independent standard Gaussians $U,V$. Hence the characteristic function of the normalised pair at $(a,b)$ converges to $e^{-(a^2+b^2)/2}$, and Lévy's continuity theorem in $\R^2$ concludes. For the argument of $\zeta$, the chain of approximations only controls $|\zeta M|$: the step $\zeta M\approx1$ in mean square does give $\arg(\zeta M)\approx0$ modulo $2\pi$, but $\arg\zeta(\frac12+it)$ is defined by continuous variation along a path, and the passage from $\sigma_0$ to $\frac12$ (Proposition~\ref{prop:offline}) has no analogue for the argument: near a zero, $\arg\zeta$ jumps by $\pi$ on a very short interval, and the average control of Lemma~\ref{lem:logint} is lost.
\end{solution}

\begin{solution}{ex:twopoint}
\idea The covariance between the two coordinates is a sum over primes that oscillates, so it stays bounded while the variances grow like $\log\log T$.
\writeup Let $\delta=2\sigma_0-1$ and $R(u)=\pi(u)-\operatorname{li}(u)$, with $\operatorname{li}(u)=\int_2^u\frac{\dd v}{\log v}$. Then
\[
\sum_{p\le X}p^{-1-\delta-i\alpha}=\int_2^X\frac{u^{-1-\delta-i\alpha}}{\log u}\dd u+\int_{2^-}^Xu^{-1-\delta-i\alpha}\dd R(u).
\]
In the first integral put $u=e^v$: it is $\int_{\log2}^{\log X}e^{-(\delta+i\alpha)v}\frac{\dd v}v$, and integrating by parts once, with $e^{-(\delta+i\alpha)v}=\frac{\dd}{\dd v}\frac{e^{-(\delta+i\alpha)v}}{-(\delta+i\alpha)}$, it is bounded by $\frac{C}{|\alpha|}$. The second, by parts, is $O(\frac1{\log^2X})+|1+\delta+i\alpha|\int_2^\infty\frac{u^{-2}u}{\log^2u}\dd u=O_\alpha(1)$. So $w_T=\sum_{p\le X}p^{-2\sigma_0-i\alpha}=O_\alpha(1)$.

Now $P_0(\sigma_0+it+i\alpha)=\sum_{p\le X}p^{-i\alpha}p^{-\sigma_0-it}$ is again a Dirichlet polynomial in $t$, with coefficients of modulus $p^{-\sigma_0}$. The proof of Proposition~\ref{prop:primesgauss} applies to any product of $j_1$ factors $P_0(s)$, $\ell_1$ factors $\overline{P_0(s)}$, $j_2$ factors $P_0(s+i\alpha)$ and $\ell_2$ factors $\overline{P_0(s+i\alpha)}$: the mean is $o(1)$ plus the diagonal sum, in which the primes of the unconjugated factors must be a permutation of those of the conjugated ones. Up to $O(v_T^{(j_1+\ell_1+j_2+\ell_2)/2-2})$ (repeated primes), the diagonal is a sum over the bijections between unconjugated and conjugated factors, each pair contributing $v_T$ if both factors belong to the same coordinate and $w_T$ or $\overline{w_T}$ otherwise. These are the moments of a complex Gaussian vector with covariance $\begin{pmatrix}v_T&w_T\\\overline{w_T}&v_T\end{pmatrix}$, and after division by $v_T$ the off-diagonal entries tend to $0$. Taking real parts and arguing as in Exercise~\ref{ex:complexgauss} (with $a\operatorname{Re}P_0(s)+b\operatorname{Re}P_0(s+i\alpha)$), the pair converges to independent standard Gaussians.
\end{solution}

\begin{solution}{ex:meandiff}
\idea The mean of $\log|\zeta(\sigma+it)|$ should not depend on $\sigma$ near $\frac12$, and the proof of Proposition~\ref{prop:offline} shows the two large contributions explicitly.
\writeup By the proof of Proposition~\ref{prop:offline}, the integral equals $\int_T^{2T}\sum_\rho f_\rho-\int_T^{2T}\bigl(\log|G(\sigma+it)|-\log|G(\frac12+it)|\bigr)\dd t$. The Gamma part is $(\sigma-\frac12)\int_T^{2T}\frac12\log\frac t{2\pi}\dd t+O(\sigma-\frac12)$, since $\operatorname{Re}\frac{G'}G(u+it)=\frac12\log\frac t{2\pi}+O(\frac1t)$. Under the Riemann hypothesis $\beta=\frac12$, so $f_\rho\ge0$ and $\int_\R f_\rho=\pi(\sigma-\frac12)$ by Lemma~\ref{lem:logint}. A zero with $\gamma\in[T,2T]$ at distance $d$ from the endpoints loses $\int_{\R\setminus[T,2T]}f_\rho\le(\sigma-\frac12)\min(\pi,\frac1d)$, and a zero outside, at distance $d$, contributes at most the same amount. There are $O(\log T)$ zeros at distance in $[j,j+1[$ from each endpoint, for $j\le T$, so these corrections total $O((\sigma-\frac12)\log T\sum_{j\le T}\frac1{j+1})+O((\sigma-\frac12)\sqrt T\log T/\sqrt T)=O((\sigma-\frac12)\log^2T)$, the zeros at distance $>T$ being handled by $|f_\rho|\le\frac{\sigma-\frac12}{d^2}$ as before. So $\int_T^{2T}\sum_\rho f_\rho=\pi(\sigma-\frac12)(N(2T)-N(T))+O((\sigma-\frac12)\log^2T)$. By Riemann--von Mangoldt, $N(2T)-N(T)=\frac1{2\pi}\bigl[u\log\frac u{2\pi e}\bigr]_T^{2T}+O(\log T)=\frac1{2\pi}\int_T^{2T}\log\frac t{2\pi}\dd t+O(\log T)$, because $\frac{\dd}{\dd u}u\log\frac u{2\pi e}=\log\frac u{2\pi}$. The main terms $\frac12(\sigma-\frac12)\int_T^{2T}\log\frac t{2\pi}\dd t$ cancel exactly.
\end{solution}

\begin{solution}{ex:fixedW}
\idea The cost of moving off the line is $(\sigma-\frac12)\log T$ on average, and Selberg's normalisation absorbs anything smaller than $\sqrt{\log\log T}$.
\writeup By Proposition~\ref{prop:offline} and Markov's inequality, $P_T\bigl(|\log|\zeta(\frac12+it)|-\log|\zeta(\sigma_T+it)||>\varepsilon\sqrt{\log\log T}\bigr)\le\frac{Cc_T}{\varepsilon\sqrt{\log\log T}}\to0$ (for $c_T=0$ there is nothing to prove). So the difference of the two normalised variables tends to $0$ in probability, and Theorem~\ref{thm:selberg} and Slutsky's lemma give the claim. For $c_T=(\log T)^{1/2}$, $\sigma_T-\frac12=(\log T)^{-1/2}$, and the Bohr--Jessen variance $\frac12\log\frac1{\sigma_T-\frac12}$ becomes $\frac14\log\log T$: half of Selberg's. Off the line by $(\log T)^{-1/2}$, only the primes up to about $e^{\sqrt{\log T}}$ contribute, and $\log\log$ of that is $\frac12\log\log T$. (The proof of this chapter can be adapted to make this heuristic rigorous; we do not pursue it.)
\end{solution}

\section*{Sources and further reading}

\begin{itemize}
\item \emph{Selberg's war.} The episodes of 1940--1945 (the fighting, the prison camp at Trandum, the arrest after the doctoral examination, the solitary research in Gj\o{}vik) and the Fields Medal of 1950 are from J.~J.~O'Connor and E.~F.~Robertson, ``Atle Selberg'', \emph{MacTutor History of Mathematics Archive}, University of St Andrews.
\item \emph{Selberg's papers.} A.~Selberg, ``On the zeros of Riemann's zeta-function'', \emph{Skr.\ Norske Vid.-Akad.\ Oslo I} 1942, no.~10. A.~Selberg, ``Contributions to the theory of the Riemann zeta-function'', \emph{Arch.\ Math.\ Naturvid.}\ 48 (1946), no.~5, 89--155. A.~Selberg, ``Old and new conjectures and results about a class of Dirichlet series'', in \emph{Proceedings of the Amalfi Conference on Analytic Number Theory} (Maiori, 1989), Univ.\ Salerno, 1992, 367--385.
\item \emph{The proof of this chapter.} M.~Radziwi\l\l{} and K.~Soundararajan, ``Selberg's central limit theorem for $\log|\zeta(\frac12+it)|$'', \emph{L'Enseignement Math.}\ (2) 63 (2017), 1--19 (arXiv:1509.06827). We follow its four propositions, its parameters and its proof of the twisted second moment, with the details written out.
\item \emph{Other proofs.} E.~Bombieri and D.~A.~Hejhal, ``On the distribution of zeros of linear combinations of Euler products'', \emph{Duke Math.\ J.}\ 80 (1995), 821--862. A.~Laurinčikas, ``A limit theorem for the Riemann zeta-function on the critical line. II'', \emph{Litovsk.\ Mat.\ Sb.}\ 27 (1987), 489--500.
\item \emph{Random matrices.} J.~P.~Keating and N.~C.~Snaith, ``Random matrix theory and $\zeta(1/2+it)$'', \emph{Comm.\ Math.\ Phys.}\ 214 (2000), 57--89.
\item \emph{Rate and tails.} A.~Roberts, ``The multivariate rate of convergence for Selberg's central limit theorem'', arXiv:2212.01411 (2022). M.~Radziwi\l\l, ``Large deviations in Selberg's central limit theorem'', arXiv:1108.5092 (2011).
\end{itemize}
